\documentclass[11pt,a4paper]{article}
\usepackage[UTF8]{ctex}
\usepackage{amsmath,amssymb,amsthm}
\usepackage{geometry}
\usepackage{tikz-cd}
\usepackage{hyperref}
\geometry{margin=2.5cm}

% Theorem environments
\newtheorem{axiom}{公理}
\newtheorem{definition}{定义}
\newtheorem{theorem}{定理}
\newtheorem{proposition}{命题}
\newtheorem{corollary}{推论}
\newtheorem{conjecture}{猜想}
\newtheorem{remark}{注}

% 定理级别标注
% [定理]：在公理1-4下有严格证明的数学命题
% [条件推论]：依赖某个猜想或特定模型条件的推论
% [诠释]：数学结构的哲学/政治经济学解释，不是数学推导
\newcommand{\lvl}[1]{\hfill\textbf{\footnotesize[#1]}}
\newcommand{\lvP}{\lvl{定理}}
\newcommand{\lvC}{\lvl{条件推论}}
\newcommand{\lvI}{\lvl{诠释}}

% Macros
\newcommand{\cat}[1]{\mathcal{#1}}
\newcommand{\Op}{\cat{O}\!p}
\newcommand{\C}{\cat{C}}
\newcommand{\I}{\mathbf{1}}
\newcommand{\lolli}{\multimap}
\newcommand{\bang}{{!}}
\newcommand{\tensor}{\otimes}
\newcommand{\coKl}[1]{\C_{\bang}}
\newcommand{\eval}{\mathrm{eval}}


\newcommand{\enc}{\mathrm{enc}}
\newcommand{\coh}{\mathrel{\frown}}
\newcommand{\eps}{\varepsilon}
\newcommand{\mU}{\mu}
\newcommand{\mN}{\nu}

\title{\textbf{操作范畴：线性逻辑、资本与生命自指}\\[6pt]
\large 生命论（明本论）的数学基础}
\author{}
\date{2026年8月}

\begin{document}
\maketitle

\begin{abstract}
本文定义\textbf{操作范畴}（operational category），将线性逻辑、范畴论、Geometry of Interaction（GoI）与马克思政治经济学统一在一个形式框架中，为生命论（明本论）提供数学基础。

\textbf{核心结果。} (1) 线性资源（活操作/活劳动）与$\bang$-模态（沉积/死劳动/资本）的结构同一：不存在$A\to\bang A$，收缩$\bang A\to\bang A\tensor\bang A$即资本增殖；(2) 反自指寄生体（资本、病毒）是$\bang$-余代数，依赖线性宿主，宿主消亡时退化为空转——资本的内在矛盾是定理；(3) 2$\times$2矩阵：交互/非交互$\times$有限/无限，$\bang$只提升非交互（数据）型，不提升交互（生命）型；(4) 不动点定理：数据函子的最大不动点可被$\bang$提升，但交互函子$F_2$的最小和最大不动点均不可——生命流$\nu F_2$不可资本化，有限交互$\mu F_2$不可复制为过程；(5) GoI动态语义：操作=带反馈的连线网络，迹=测量=实际从属，三者范畴论同一；(6) 余归纳GoI与生产性迹范畴（PTC）：以productivity替换nilpotency，公理化无限生产性过程，$\bang$不穿透任意PTC中的线性状态；(7) f-层级：自我模型必然是$\bang$-模态的（意识形态的存在论基础），异化=错认$\bang$-模型为线性过程，明性=看穿模型为模型，f$^3$是层级提升的不动点（$\alpha\in\{1,2,3\}$）；(8) 可靠性与完全性：可证的生命型过程解释为生产性流处理器，但生命过程超越有限证明（生产性流处理器不可数，证明可数）；(9) 量子对应：不可克隆定理是$\bang$-不穿透定理在Hilbert空间模型中的特例，量子测量=GoI迹=实际从属；(10) $G^\omega(\mathcal{C})$构造：PTC上的dagger紧闭合范畴，递归类型可解释，$\bang$限制仍成立；(11) 量论操作化：$T$=幂零指数，$N$=谱半径，$\alpha$=反馈嵌套深度，$M=\alpha TN$，慢性死亡定理（$N<1$时$M$有限）；(12) 革命=消除$\bang$-寄生Cut，社会革命与个人明性是同一操作；(13) 生命论基本定理：九条命题共享同一数学结构——线性资源不可被!-模态还原。其中四条（第2、3、7、9条）有严格证明，为[定理]；两条（第4、5条）的范畴论区分本身是[定理]，与"异化""明性"的等同是[诠释]；三条（第1、6、8条）为[诠释]。"共享同一结构"指：每条命题都能从种子定理1（不存在A→!A）或其经不动点/GoI结构的推广（定理14、20、38）与线性逻辑公理的一个小推导中得到。

本文的贡献包括：(1) 将这些组件组装为操作范畴并赋予生命论解释；(2) 证明了关于$\bang$-模态与线性资源的若干定理；(3) 给出了马克思政治经济学核心概念的范畴论形式化；(4) 提出了践演判断（performative judgment）作为新的判断形式。
\end{abstract}

\section{操作范畴的定义}

\begin{definition}[操作范畴]\label{def:opcat}
一个\textbf{操作范畴}是六元组
$$\Op = (\C, \tensor, \lolli, \I, \bang, \mU/\mN)$$
满足以下公理。
\end{definition}

\begin{axiom}[闭范畴性]\label{ax:linear}
$(\C, \tensor, \I)$是对称幺半范畴，且$\C$是幺半闭的：对任意对象$A,B$，存在内部hom $A\lolli B$与求值态射
$$\eval_{A,B} : (A\lolli B)\tensor A \to B,$$
满足curry--de-curry同构$\C(A\tensor B,C)\cong\C(A,B\lolli C)$。

对象称为\textbf{操作位点}（operational site）；$A\tensor B$表示两个独立操作并行；$\I$是空操作；$A\lolli B$表示消耗$A$的一次运行产生$B$的一次运行。
\end{axiom}

\begin{axiom}[无自然收缩]\label{ax:no-contraction}
不存在自然变换$\Delta_A : A \to A\tensor A$（对角）与$p_A : A \to \I$（投影）。态射无自复制，线性资源不可任意复制或丢弃。
\end{axiom}

\begin{axiom}[指数余单子]\label{ax:bang}
$(\bang, \eps, \delta)$是$\C$上的\textbf{线性指数余单子（linear exponential comonad）}，满足：
\begin{itemize}
  \item (i) 余单子结构：$\eps_A : \bang A \to A$（弃置）、$\delta_A : \bang A \to \bang\bang A$（升进）；
  \item (ii) 余交换余幺半群结构：$\kappa_A : \bang A \to \bang A\tensor\bang A$（复制）、$\omega_A : \bang A \to \I$（弱化），即每个$\bang A$带余交换余幺半群结构；
  \item (iii) \textbf{Seely同构}：$\bang(A\&B)\cong\bang A\tensor\bang B$且$\bang\top\cong\I$（其中$\&$为笛卡尔积、$\top$为终对象）。
\end{itemize}
注（[诠释]）：上述结构在生命论中被称为"阴化余单子"；此命名为[诠释]，不改变公理的数学内容。
\end{axiom}

\begin{axiom}[不动点]\label{ax:fixedpoints}
对本文使用的\textbf{线性公式函子}——由$\tensor,\&,\oplus,\lolli$连接词在变元$X$正出现位置生成的函子$F(X)$——存在最小不动点$\mU F$（初始代数）与最大不动点$\mN F$（最终余代数），以及规范态射$\iota:\mU F\to\mN F$。
\end{axiom}

\noindent\textbf{公理化的一致性。}公理1--4非空：$\mathbf{Rel}$（关系范畴，以集合的笛卡尔积为$\tensor$、以$\bang A$为$A$上的有限多重集）满足公理1--4（含Seely条件），且本文所用线性公式函子在$\mathbf{Rel}$中有$\mu/\nu$不动点。加权$\mathbf{Rel}$进一步为量论参数提供模型（§20）。$\mathbf{Hilb}^\omega$与$\mathbf{Coh}^\omega$为候选模型（§22，严格验证见猜想）。

\begin{remark}[践演元公理——哲学预设，非数学公理]\label{ax:performative}
$\Op$的任何使用（构造态射、画图、追图、执行求值）本身是$\Op$之外的一次活操作。不存在元范畴能完全表示$\Op$的所有活操作。

\textbf{状态说明}：这不是$\C$内部的数学公理，而是关于$\C$的哲学预设。它不可形式化（任何形式化本身都是一次表示，而该预设断言表示不能穷尽活操作）。本文的数学定理（标注[定理]）不依赖此预设；标注[诠释]的命题以此预设为前提。
\end{remark}

\section{核心定理}

\begin{theorem}[阳不可自动阴化]\label{thm:no-A-to-bangA}\lvP
在操作范畴中，不存在自然变换$\eta_A : A \to \bang A$。
\end{theorem}

\begin{proof}
假设存在自然变换$\eta_A : A \to \bang A$。考虑复合
$$A \xrightarrow{\eta_A} \bang A \xrightarrow{\kappa_A} \bang A\tensor\bang A \xrightarrow{\eps_A\tensor\eps_A} A\tensor A.$$
这给出自然变换$A\to A\tensor A$，即对角化$\Delta_A$，与公理\ref{ax:no-contraction}矛盾。
\end{proof}

\begin{remark}
定理\ref{thm:no-A-to-bangA}说：活操作不能自动变成可复制的沉积。记录、表达、对象化需要额外的沉积工作——promotion规则$\bang\Gamma\vdash A \;/\; \bang\Gamma\vdash\bang A$要求已有沉积$\bang\Gamma$作为支撑。
\end{remark}

\begin{remark}[种子定理]\label{rmk:seed}
定理\ref{thm:no-A-to-bangA}——不存在自然变换$A\to\bang A$——是本文的\textbf{种子定理}。本文"生命不可资本化"系列定理（6、16、17、18、20、23、38）均为其直接实例（取特定$A$或$S$）或经不动点/GoI结构的派生。这一事实构成定理51第9条声称的严格内容。
\end{remark}

\begin{theorem}[阴不能穷尽阳]\label{thm:yin-cannot-exhaust-yang}\lvP
在操作范畴中，Lawvere不动点定理的证明无法进行：该证明需要对角化$\Delta : A\to A\tensor A$来构造自我应用$\phi(a)(a)$，但公理\ref{ax:no-contraction}禁止对角化。沉积$\bang A$可以通过$\kappa$复制再经$\eps$弃置得到$\bang A \to A\tensor A$，但这给出的是$\bang A\to A\tensor A$而非$A\to A\tensor A$——阴可以模拟阳的复制，但模拟不等于原初运行。
\end{theorem}

\begin{theorem}[最终余代数的存在性]\label{thm:productive-self-ref}\lvP
设$\mN F$是$F$的最终余代数。对任何余代数$f:X\to F(X)$，存在唯一的余代数同态$\mathrm{beh}_f : X\to\mN F$使得
$$\mathrm{out}\circ\mathrm{beh}_f = F(\mathrm{beh}_f)\circ f.$$
\end{theorem}

\begin{proof}
这是最终余代数的定义性性质（finality）。
\end{proof}

\begin{remark}
定理\ref{thm:productive-self-ref}说：活自指（展开$S\to F(S)$，而非自我应用$\phi(a)(a)$）总是良基的，有唯一解，不导致悖论。生命的自指是展开性的（今天$\to$明天），不是对角化的（把自己复制一份来否定自己）。
\end{remark}

\section{反自指寄生体}

\begin{definition}[反自指寄生体]\label{def:parasite}
设$S=\mN F$是活自指系统（最终余代数）。一个\textbf{反自指寄生体}$P$是$\bang$-模态对象$P=\bang Q$，配有：
\begin{enumerate}
  \item \textbf{抽取}：$d : P\tensor S \to P\tensor P$（$P$消耗$S$复制自身）；
  \item \textbf{消耗}：$c : P\tensor S \to S'$，其中不存在$s':S'\to S$（削弱不可逆）；
  \item \textbf{依赖性}：$P$不具有独立的余代数结构$P\to F(P)$。
\end{enumerate}
\end{definition}

\begin{theorem}[宿主削弱至幺元时抽取态射退化为纯收缩]\label{thm:parasite-self-destruct}\lvP
若宿主状态$S$被消耗削弱至幺元$\I$，则复合态射$d:P\tensor S\to P\tensor P$退化为纯收缩$P\to P\tensor P$。
\end{theorem}

\begin{proof}
$S=\I$时$d:P\tensor\I\to P\tensor P$，由张量单位律即$P\to P\tensor P$；无额外结构可阻止。\qed
\end{proof}

\begin{remark}[空转——条件推论]\lvC
在GoI执行语义下（沉积的激活须经一次执行），$S=\I$无可执行的活操作，$\eps:\bang Q\to Q$无法被实际调用，$P$的扩张仅是沉积层面的数字膨胀，无任何运行被激活。此结论依赖执行语义，非纯范畴论推论。
\end{remark}

\begin{remark}[马克思"资本的绝对矛盾"——诠释]\lvI
上式被解释为马克思"资本的绝对矛盾"：寄生体消灭宿主即消灭自身的增殖条件。此等同不作数学断言。
\end{remark}

\section{商品与资本}

\begin{definition}[价值量，模型内]\label{def:value-magnitude}
对$\bang$-模态对象$\bang C$，定义其价值量$V(\bang C)$为$\bang C$在余代数结构（counit/comultiplication）下的同构不变量，即$V(\bang C\tensor\bang C)\cong V(\bang C)$。

注：将此$V$命名为"价值量"是[诠释]；"$V(\bang C\tensor\bang C)\cong V(\bang C)$"本身是comonoid律的纯数学推论，为[定理]。
\end{definition}

\begin{definition}[商品]\label{def:commodity}
一个\textbf{商品}是对象$A\tensor\bang V$，其中$A$（线性组件）是\textbf{使用价值}（具体的、在使用中消耗的），$\bang V$（$\bang$-模态组件）是\textbf{交换价值}（抽象的、可复制可积累的）。
\end{definition}

\begin{definition}[货币与资本]\label{def:money-capital}
\textbf{货币}是$\bang G$——纯交换价值，一般等价物。\textbf{资本}是共Kleisli范畴中的态射$k:G\to G'$（对应$\C$中的$\bang G\to\bang G'$），满足$G'>G$（增殖）。在基范畴中，资本流通需要线性中介$A$（劳动力）：
$$p : \bang C\tensor A \to \bang C\tensor B.$$
\end{definition}

\begin{theorem}[死劳动/资本的保存与活化]\label{thm:surplus-linear}\lvC
在操作范畴中：
\begin{enumerate}
  \item[(a)] 由comonoid律，$V(\bang C\tensor\bang C)\cong V(\bang C)$，价值量在收缩之下不变（严格说是值域同构，价值转移而非创造，因为$\Delta_A$不存在、$\bang C$自身不带投射）；
  \item[(b)] 若$X\in\mathrm{Ob}(\C)$线性、无收缩（$X\nrightarrow X\tensor X$），则$X\tensor\bang C\to Y\tensor\bang C$中$X$被消耗，$C$的$\bang$-模式无法将$X$复制保存；
  \item[(c)] $\bang C\to C$（$\eps$）不增加价值量——它只是激活已沉积的量，不产生新量；
  \item[(d)] 收缩是活动的源泉：任何真实增殖须经$\bang C$的$\kappa$，即资本增殖的本质是沉积结构的自我复制，而非线性资源的再生。
\end{enumerate}
注（[诠释]）：(a)(c)的"价值量"依定义\ref{def:value-magnitude}；(b)(d)中的"资本"即$\bang$-模态增殖结构。上述陈述中，除定义内的同构推论为[定理]外，其余为[条件推论]（依赖抽象"价值量"的设定）。
\end{theorem}

\begin{theorem}[最大不动点对象无到!-模态的态射]\label{thm:life-not-capital}\lvP
设$S=\mN F$为生命流主体（活过程），不存在态射$h:S\to\bang S$。
\end{theorem}

\begin{proof}
取定理\ref{thm:no-A-to-bangA}之$A=S$，直接得到不存在$A\to\bang A$。\qed
\end{proof}

\begin{remark}[异化的数学限度——诠释]\lvI
异化有数学限度：生命过程不可被完全沉积为可复制的$\bang$-模态结构。这里的"异化"是下述结构在生命论中的[诠释]命名——不存在$S\to\bang S$这一数学事实本身为[定理]。
\end{remark}

\section{流通领域与生产领域}

\begin{theorem}[Benton 1995]\label{thm:cokleisli-ccc}\lvP
设$(\bang,\eps,\delta)$是SMCC $\C$上的线性指数余单子，则其共Kleisli范畴$\coKl{!}$是笛卡尔闭范畴（CCC）。
\end{theorem}

\begin{theorem}[共Kleisli范畴的结构]\label{thm:fetishism}\lvP
遗忘函子$U:\C\to\coKl{!}$将线性态射$f:A\to B$映射为共Kleisli态射$\bang A\to B$；在$\coKl{!}$中，$A$与$\bang C$均为交换幺半群对象（具对角化与投影）。因此，在$\coKl{!}$中不存在范畴论性质能够内在地区分$A$与$\bang C$。
\end{theorem}

\begin{proof}
$U$的结构如定义；在$\coKl{!}$中$A$是普通对象（可通过$\bang$提升为$\bang A$），$\bang C$是余代数对象，两者都在$\coKl{!}$中成为可投影/可复制的幺半群对象。\qed
\end{proof}

\begin{remark}[商品拜物教——诠释]\lvI
上述结构被解释为马克思的商品拜物教：$\C$（生产领域）中$A$（线性活劳动）与$\bang C$（$\bang$-模态资本）有结构区别、剥削可见；$\coKl{!}$（流通领域）中一切皆为$\bang$-模态、剩余价值来源不可见。"$X$对应$Y$"不声称"$X$就是$Y$"。
\end{remark}

\begin{theorem}[沉积积累与线性资源的增长不对称]\label{thm:organic-composition}\lvC
设生产态射序列$p_n:\bang C_n\tensor A_n\to\bang C_{n+1}$表示$n$期生产。则$\bang C_n$可通过收缩无结构障碍地积累，而$A_n$是线性的、不能收缩，其增长只能外延扩大。在剩余价值仅来自$A$的供给模型中（定理\ref{thm:surplus-linear}条件），$\bang C_n$的积累快于$A_n$的增长。
\end{theorem}

\begin{proof}
每期产出$\bang C_{n+1}$包含剩余$\bang G_m$，通过promotion沉积为新的不变资本。$\bang C$是$\bang$-模态，积累通过$\kappa$无结构障碍。$A$是线性的，公理\ref{ax:no-contraction}禁止$A\to A\tensor A$，增加$A$只能外延扩大。
\end{proof}

\begin{remark}[有机构成上升与利润率下降——诠释]\lvI
上述不对称在政治经济学中被解释为马克思"资本有机构成上升"与"利润率趋向下降"。"利润率"须在具体模型中另行定义；此等同不作数学断言。
\end{remark}

\section{践演坐实}

\begin{definition}[践演判断]\label{def:performative-judgment}
除标准判断$\Gamma\vdash A$（消耗$\Gamma$产生$A$）外，引入\textbf{践演判断}
$$\triangleright A,$$
读作"$A$在操作上坐实"。$\triangleright A$不是类型（不是$\C$中的对象），不能出现在项的位置。
\end{definition}

\begin{theorem}[运行真理]\label{thm:operational-truth}\lvI
设$L$是足够强的形式系统（其证明论可表示为$\bang$-模态对象）。设$\mathrm{Op}_L$为"$L$在运行"。则：
\begin{enumerate}
  \item $\mathrm{Op}_L$不能在$L$中证明（它不是$L$中的命题，而是$L$运行的元条件）；
  \item $\mathrm{Op}_L$不能在$L$中否定（否定$\mathrm{Op}_L$的证明本身是$L$中的一次操作，线性消耗$\mathrm{Op}_L$后矛盾不传播）；
  \item $L$的一致性$\mathrm{Con}(L)$是$\mathrm{Op}_L$的$\bang$-模态沉积；
  \item $\mathrm{Con}(L)$不可在$L$中证明（G\"odel第二定理），这是定理\ref{thm:no-A-to-bangA}的特例。
\end{enumerate}
\end{theorem}

\section{生命四规定性}

\begin{definition}[生命位点]\label{def:life-site}
设$F(S)=(S\tensor O)^I$为Mealy机函子。$F$-余代数$(S,f:S\to F(S))$是\textbf{生命位点}，如果：
\begin{enumerate}
  \item \textbf{边界}：$S$是明确对象，区分内部（状态$S$）与外部（输入$I$、输出$O$）；
  \item \textbf{内生目的}：存在存活判据$v:S\to 2$与反馈$m:O\to S$，使输出的一部分反馈维持$S$；
  \item \textbf{再生}：$\mathrm{out}:S\to F(S)$每步产生新的$S$（$S$在$F(S)$中正位置出现）；
  \item \textbf{互动}：$I$和$O$非平凡。
\end{enumerate}
\end{definition}

\begin{definition}[f-层级]\label{def:f-levels}
生命位点按余代数嵌套深度分层：
\begin{itemize}
  \item $f^1$（自在）：$S_1\to(S_1\tensor O_1)^{I_1}$（简单自维持）；
  \item $f^2$（自为）：$S_2\to(S_2\tensor S_2\tensor O_2)^{I_2}$（含内部模型）；
  \item $f^3$（自觉/明性）：$S_3\to(S_3^{S_3}\tensor O_3)^{I_3}$（含自我模型$S_3\lolli S_3$）。
\end{itemize}
\end{definition}

\begin{proposition}[f-层级与寄生条件]\label{prop:f3-alienation}\lvI
反自指寄生体$P=\bang Q$在不同f-层级的作用方式有本质区别：
\begin{itemize}
  \item $f^1$：无自我模型，$\bang$-模态无劫持对象——寄生体只能通过外部物理强制，不能通过意识形态（模型替换）运作。
  \item $f^2$：有$\bang$-模态自我模型$\bang\ulcorner f\urcorner$（定理\ref{thm:self-model-bang}），但不能区分模型与运行。寄生体替换模型后$f^2$直接按新模型运行——这是\textbf{被控制}，不是\textbf{被异化}：$f^2$没有"选择是否认同模型"的能力，它不知道自己被劫持了。
  \item $f^3$：能区分模型与运行（定理\ref{thm:mingxing}）。寄生体需要$f^3$\emph{自愿}把$\bang$-模型认同为自己（把沉积模型错认为线性运行，即定理\ref{thm:alienation-mistake}的异化条件）才能劫持——这才是\textbf{异化}：异化预设了选择能力，也只有有选择能力的存在才能被异化。$f^3$既能被异化，也能看穿模型而明性。
\end{itemize}

数学内核：$f^1$状态无$\bang$-组件（[定理]）；$f^2$有$\bang$-模型但转移函数无"判断"步骤（[定义]）；$f^3$有$\bang H_3$且转移函数含模型-运行区分（定理\ref{thm:f3-fixedpoint}、\ref{thm:mingxing}）。"寄生""控制""异化"的命名为[诠释]。
\end{proposition}

\section{交互余代数不可资本化}

\begin{proposition}[$\bang$提升到非交互余代数]\label{prop:lift-noninteractive}\lvP
对$F_1(X)=X\tensor B$（无输入的自复制），分配律$\lambda^1:\bang F_1\Rightarrow F_1\bang$存在。
\end{proposition}
\begin{proof}
由$\bang$的monoidal strength，$\bang(X\tensor B)\to\bang X\tensor\bang B$，再与$\eps_B:\bang B\to B$复合，得到$\bang(X\tensor B)\to\bang X\tensor B=F_1(\bang X)$。
\end{proof}

\begin{theorem}[$\bang$不提升到交互余代数]\label{thm:no-lift-interactive}\lvP
对$F_2(X)=A\lolli(X\tensor B)$（Mealy机函子，交互自复制），其中$A,B$为非$\bang$-模态对象，分配律$\lambda^2:\bang F_2\Rightarrow F_2\bang$不存在。即矢列
$$\bang(A\lolli(X\tensor B)),\, A\ \vdash\ \bang X\tensor B$$
在直觉主义乘法指数线性逻辑中不可证。
\end{theorem}

\begin{proof}
结论$\bang X\tensor B$的主连接词是$\tensor$，最后一步只能是$\tensor R$：
$$\frac{\Gamma_1\vdash\bang X\qquad \Gamma_2\vdash B}{\Gamma_1,\Gamma_2\vdash\bang X\tensor B}$$
其中$\Gamma_1\cup\Gamma_2=\{\bang(A\lolli(X\tensor B)),\,A\}$。四种分配：
\begin{enumerate}
  \item $\Gamma_1=\emptyset$：$\vdash\bang X$要求promotion给出$\vdash X$，空上下文不可证原子$X$。
  \item $\Gamma_2=\emptyset$：$\vdash B$不可证。
  \item $\Gamma_1=\{\bang(A\lolli(X\tensor B))\},\ \Gamma_2=\{A\}$：$A\vdash B$不可证（不同原子）。
  \item $\Gamma_1=\{A\},\ \Gamma_2=\{\bang(A\lolli(X\tensor B))\}$：$A\vdash\bang X$要求promotion，但$A$不是$\bang$-模态，不适用。
\end{enumerate}
四种分配均不可证。结论不以$\bang$开头，promotion不是最后一步。因此矢列不可证。
\end{proof}

\begin{remark}
即使输入也是$\bang$-模态（$\bang A$），矢列仍不可证：弃置$\bang A$得$A$，但$A\vdash X$不可证。问题在输出端——函数产生的$X$是线性的，无法变成$\bang X$。
\end{remark}

\begin{theorem}[交互过程不可被!-提升为复制的交互]\label{thm:subsumption}\lvP
由命题\ref{prop:lift-noninteractive}与定理\ref{thm:no-lift-interactive}，存在分配律$\bang F_1\Rightarrow F_1\bang$，不存在分配律$\bang F_2\Rightarrow F_2\bang$。因此：将交互式过程改造为非交互式（$F_2\to F_1$）后其可被$\bang$提升为复制品；保持交互式的过程不可被整体复制。
\end{theorem}

\begin{remark}[形式从属与实际从属——诠释]\lvI
上式被解释为马克思"形式从属不可能，实际从属是资本增殖的唯一途径"：资本（$\bang$-模态增殖结构）无法以复制方式吸收活交互（$F_2$），只能先把活交互改造为死流程（$F_2\to F_1$）再行支配。"资本"与"实际从属"的等同不作数学断言。
\end{remark}

\section{数据与计算：!-分配的一般刻画}

我们将定理\ref{thm:no-lift-interactive}推广到一般线性公式函子，用四个递归判断精确刻画。

\begin{definition}[四个提升判断]
对含自由变元$X$的公式$G(X)$，定义：
\begin{itemize}
  \item $P(G)$：$G(X)\vdash G(\bang X)$可证（线性$G$可直接提升）
  \item $D(G)$：$G(\bang X)\vdash G(X)$可证（$\bang$-版本可弃置回线性）
  \item $Q(G)$：$\bang G(X)\vdash G(\bang X)$可证（\textbf{分配律}，即$\bang$推过$G$）
  \item $R(G)$：$\bang G(\bang X)\vdash G(X)$可证（双$\bang$可弃置回线性）
\end{itemize}
分配律$\lambda^G:\bang G(X)\to G(\bang X)$存在当且仅当$Q(G)$成立。
\end{definition}

\begin{theorem}[递归刻画]\label{thm:recursive-characterization}\lvP
四个判断按以下规则递归计算：

\noindent 基例：$P(X)=\bot$，$D(X)=Q(X)=R(X)=\top$；对不含$X$的常量$C$，四个判断均为$\top$。

\noindent 乘法$\tensor$：$P,D,Q,R$均按分量合取（如$Q(G_1\tensor G_2)=Q(G_1)\land Q(G_2)$），由monoidal strength$\bang(G_1\tensor G_2)\to\bang G_1\tensor\bang G_2$给出。

\noindent 加法与$\&$：与$\tensor$相同（由Seely同构，公理3iii）。

\noindent 加法析取$\oplus$：$P,D,R$按分量合取，但$Q(G_1\oplus G_2)=\bot$（当$X$在任一分支中正出现）。

\noindent 线性蕴含$\lolli$：
\begin{align*}
P(G_1\lolli G_2) &= D(G_1)\land P(G_2)\\
D(G_1\lolli G_2) &= P(G_1)\land D(G_2)\\
Q(G_1\lolli G_2) &= D(G_1)\land P(G_2)\\
R(G_1\lolli G_2) &= P(G_1)\land D(G_2)
\end{align*}

\noindent 指数$\bang$：$P(\bang G_1)=Q(G_1)$，$D(\bang G_1)=R(G_1)$，$Q(\bang G_1)=Q(G_1)$，$R(\bang G_1)=R(G_1)$。
\end{theorem}

\begin{proof}
$\oplus$情形：$\bang(G_1\oplus G_2)\vdash G_1(\bang X)\oplus G_2(\bang X)$的最后一步只能是$\oplus R_1$或$\oplus R_2$。$\oplus R_1$要求$\bang(G_1\oplus G_2)\vdash G_1(\bang X)$，弃置后$\oplus L$要求$G_1\vdash G_1(\bang X)$（即$P(G_1)$）和$G_2\vdash G_1(\bang X)$。若$X$在$G_1$中正出现，$P(G_1)$失败（需要$X\to\bang X$，由定理\ref{thm:no-A-to-bangA}不可能）。$\oplus R_2$同理。promotion不能作为最后一步（结论不以$\bang$开头）。

$\lolli$情形：以$P(G_1\lolli G_2)$为例，需证$(G_1\lolli G_2)\vdash G_1(\bang X)\lolli G_2(\bang X)$，Curry反转后为$(G_1\lolli G_2),G_1(\bang X)\vdash G_2(\bang X)$。用$D(G_1)$将$G_1(\bang X)$降级为$G_1(X)$，应用$\lolli L$得$G_2(X)$，再用$P(G_2)$提升为$G_2(\bang X)$。$Q(G_1\lolli G_2)$相同（弃置$\bang(G_1\lolli G_2)$后与$P$情形一致）。$D$和$R$对偶。

$\bang$情形：$P(\bang G_1)$由promotion归约为$Q(G_1)$；$D(\bang G_1)$由promotion归约为$R(G_1)$；$Q(\bang G_1)$弃置$\bang\bang G_1$后用$Q(G_1)$；$R(\bang G_1)$弃置后用$R(G_1)$。
\end{proof}

\begin{theorem}[!-分配二分法]\label{thm:dichotomy}\lvP
对任意线性公式函子$F$，分配律$\bang F(X)\to F(\bang X)$存在当且仅当$Q(F)=\top$。具体地，$Q(F)$失败当且仅当$X$出现在$\lolli$的\emph{后件}（codomain，被函数生产出来的位置）或$\oplus$的分支中。
\end{theorem}

\begin{proof}
由定理\ref{thm:recursive-characterization}对$F$结构归纳。关键情形：$F(X)=A\lolli G(X)$时$Q(F)=D(A)\land P(G)$（$A$不含$X$故$D(A)=\top$），而$P(G)$在$X$正出现于$G$时为$\bot$。$F(X)=F_1\oplus F_2$时$Q(F)=\bot$。其余情形$Q$均成立。
\end{proof}

\begin{remark}[术语]
"线性公式函子"指由线性连接词$\tensor,\&,\oplus,\lolli$在变元正出现位置生成的函子；$\lolli$在逆变位置不是协变函子，故不属于一般多项式函子理论，此处命名不承诺多项式函子文献中的任何结论。
\end{remark}

\begin{remark}[修正]
早期版本声称"$X$在$G$中任意出现即阻止$A\lolli G$的分配律"，这是错误的。反例：$F(X)=X\lolli B$（$X$在domain/输入位置），$Q(F)=D(X)\land P(B)=\top\land\top=\top$，分配律$\bang(X\lolli B)\to\bang X\lolli B$存在（弃置后与$\bang X$弃置为$X$求值得$B$）。只有$X$在\emph{codomain}（被生产出来）时才阻止。生命$F(X)=A\lolli(X\tensor B)$中$X$在codomain，定理\ref{thm:no-lift-interactive}和\ref{thm:subsumption}的结论不受影响。
\end{remark}

\begin{remark}[与严格正性的对偶]
定理\ref{thm:dichotomy}中"$X$不出现在$\lolli$后件中"与归纳类型的严格正性条件形成对偶：严格正性禁止$X$出现在$\to$的\emph{domain}（否则非良基递归），而!-分配禁止$X$出现在$\lolli$的\emph{codomain}（否则计算结果不可复制）。阴（归纳/$\mu$）禁止$X$作为输入；阳（!-分配/资本化）禁止$X$作为输出。
\end{remark}

\begin{remark}[与量子no-go定理的统一]
线性逻辑是量子计算的类型论基础。定理\ref{thm:dichotomy}将量子信息中三个独立的no-go定理统一为$Q(F)$的不同子句：$P(X)=\bot$对应不可克隆定理（Wootters-Zurek 1982）；$Q(A\lolli(X\tensor B))=\bot$对应不可编程定理（Nielsen-Chuang 1997）；$Q(X\oplus B)=\bot$对应不可广播定理（Barnum et al. 1996）。生命作为交互计算过程$A\lolli(X\tensor B)$，与量子过程一样不可复制。
\end{remark}

\section{不动点的!-分配：生命流不可资本化}

我们将定理\ref{thm:dichotomy}推广到不动点。设$G(X)=F(\bang X)$，问$\bang\mu F\to\mu G$和$\bang\nu F\to\nu G$是否存在。

\begin{theorem}[数据函子的$\nu$-提升]\label{thm:nu-lift-data}\lvP
若$Q(F)=\top$（分配律$\lambda^F:\bang F\to F\bang$存在），则$\bang\nu F\to\nu(F\bang)$存在。
\end{theorem}

\begin{proof}
设$\mathrm{out}_F:\nu F\to F(\nu F)$为最终余代数。构造$\bang\nu F$上的$G$-余代数（$G(X)=F(\bang X)$）：
$$\bang\nu F\xrightarrow{\delta}\bang\bang\nu F\xrightarrow{\bang\,\mathrm{out}_F}\bang F(\nu F)\xrightarrow{\lambda^F_{\nu F}}F(\bang\nu F)\xrightarrow{F(\delta)}F(\bang\bang\nu F)=G(\bang\nu F).$$
由$\nu G$的终结性，存在唯一态射$\bang\nu F\to\nu G$。
\end{proof}

\begin{theorem}[交互函子最大不动点的!-分配不存在]\label{thm:nu-no-lift-life}\lvP
对$F_2(X)=A\lolli(X\tensor B)$，不存在$\bang\nu F_2\to\nu(F_2\bang)$。
\end{theorem}

\begin{proof}
需证$\bang\nu F_2\vdash\nu G$，$G(X)=A\lolli(\bang X\tensor B)$。由$\nu$-fold，只需证$\bang\nu F_2\vdash A\lolli(\bang\nu G\tensor B)$，即$\bang\nu F_2,A\vdash\bang\nu G\tensor B$。结论主连接词为$\tensor$，最后一步只能是$\tensor R$，四种上下文拆分与定理\ref{thm:no-lift-interactive}完全相同：(1)空上下文证$\bang\nu G$需promotion给$\vdash\nu G$，再由$\nu$-fold要求$A\vdash\bang\nu G\tensor B$，$A$为新原子不可证；(2)$\vdash B$不可证；(3)$A\vdash B$不可证；(4)$A\vdash\bang\nu G$要求promotion但$A$非$\bang$-模态。左侧收缩、弃置、$\nu$-unfold不改变结论：弃置后unfold求值得$\nu F_2\tensor B$，拆分后仍需$\nu F_2\vdash\bang\nu G$（promotion不适用）或回到同一目标（循环）。

\textbf{方法论说明}：不可证性的完整论证采用证明搜索完备性（焦点化，如Andreoli的LKF），只需穷举主规则。四种$\tensor R$拆分是结论主公式的全部可能；左侧$\nu$-unfold/$\mu$-elimination不产生新的可证分支（unfold后目标复杂度不降低，归纳步骤与弃置后目标相同，构成循环而非进展）。
\end{proof}

\begin{theorem}[有限交互过程不可提升为过程]\label{thm:mu-no-lift}\lvP
对$F_2(X)=A\lolli(X\tensor B)$，不存在$\bang\mu F_2\to\mu(F_2\bang)$。
\end{theorem}

\begin{proof}
由$\mu$-fold，需证$\bang\mu F_2,A\vdash\bang\mu G\tensor B$，与定理\ref{thm:nu-no-lift-life}结构相同（$\nu$换$\mu$），四种$\tensor R$拆分同样全部失败。$\mu$-elimination归纳步骤$F_2(\bang\mu G\tensor B)\vdash\bang\mu G\tensor B$经$\lolli L$分解后，一个子目标成立（$A\vdash A$公理），另一个子目标成立（弃置多余$B$），但归纳结论$\bang\mu F_2,A,\mu F_2\vdash\bang\mu G\tensor B$与弃置后目标相同，构成循环而非进展。
\end{proof}

\begin{remark}[种子定理派生的两个方向]
种子定理\ref{thm:no-A-to-bangA}（不存在$A\to\bang A$）经不动点结构有两类推广：\\
(i) \textbf{对象级}：取$S=\nu F_2$（定理\ref{thm:life-not-capital}）或$A=\mu F_2$（定理\ref{thm:no-A-to-bangA}），得$\nu F_2\to\bang\nu F_2$、$\mu F_2\to\bang\mu F_2$不存在——阴化方向被阻；\\
(ii) \textbf{结构级}：定理\ref{thm:nu-no-lift-life}、\ref{thm:mu-no-lift}得$\bang\nu F_2\to\nu(F_2\bang)$、$\bang\mu F_2\to\mu(F_2\bang)$不存在——$\bang$-分配方向被阻。\\
两个方向彼此独立，共同构成"生命流不可资本化"的两翼。
\end{remark}

\begin{remark}[2$\times$2矩阵]
分界不是有限/无限，而是交互/非交互：
\begin{center}
\begin{tabular}{c|cc}
 & $\mu$（有限） & $\nu$（无限）\\\hline
$F_1$（非交互/数据） & \checkmark & \checkmark（定理\ref{thm:nu-lift-data}）\\
$F_2$（交互/生命） & $\times$（定理\ref{thm:mu-no-lift}） & $\times$（定理\ref{thm:nu-no-lift-life}）
\end{tabular}
\end{center}
数据（无论有限无限）可被$\bang$复制；交互过程（无论有限无限）不可被$\bang$复制为过程。有限交互可被\emph{记录为数据}（转化为$F_1$型结构），但记录不是交互——这正是定理\ref{thm:subsumption}的实际从属。

\textbf{内层与外层$\bang$。}本文同时证明两个方向的$\bang$不穿透：(1) \textbf{内层}$\bang\nu F_2\to\nu(F_2\bang)$（$\bang$进入每个展开步，定理\ref{thm:nu-no-lift-life}、\ref{thm:mu-no-lift}）；(2) \textbf{外层}$\bang\nu F_2\to\nu\bang F_2$（整个不动点被$\bang$包裹，定理\ref{thm:goi-no-clone}、\ref{thm:mu-rigorous}）。二者共享同一本质（反馈线$X$的线性性），但作为数学命题彼此独立。
\end{remark}

\begin{remark}[哲学意义]
定理\ref{thm:nu-no-lift-life}是"生命不可完全资本化"（定理\ref{thm:life-not-capital}）的加强：不仅一步活操作不可资本化，整个无限交互自指过程$\nu F_2$——一个持续回应、持续成长的生命——不可被$\bang$穿透。这是宇宙中唯一不可被复制/沉积/资本化的结构：无限的交互自指过程。资本唯一能做的是把交互改造成非交互（去技能化，定理\ref{thm:subsumption}），把生命过程改造成数据流。
\end{remark}

\section{GoI动态语义：操作的执行}

范畴论语义中态射$f:A\to B$是静态的。Girard的Geometry of Interaction（GoI）给出动态语义：态射是带反馈回路的连线（wiring），cut是反馈，执行是反馈的迭代（Girard 1989; Abramsky 1996用traced monoidal范畴形式化）。

\begin{definition}[迹]\label{def:trace}
在traced symmetric monoidal范畴中，迹算子
$$\mathrm{Tr}^X_{A,B}(f):A\to B,\qquad f:A\tensor X\to B\tensor X$$
把输出$X$接回输入$X$形成反馈回路，图形上表示为：
$$A\longrightarrow\boxed{f}\longrightarrow B,\qquad X\longrightarrow\boxed{f}\longrightarrow X\ \text{（反馈）}.$$
\end{definition}

\begin{proposition}[Mealy机是GoI反馈回路]\label{prop:mealy-goi}\lvP
$F_2$-余代数$f:A\tensor X\to B\tensor X$（Mealy机）在GoI中是带反馈的盒子：$A$为外部输入，$B$为外部输出，$X$为内部状态（自我），输出$X$接回输入$X$。迹$\mathrm{Tr}^X_{A,B}(f):A\to B$藏起状态，只留可观察的输入输出行为。
\end{proposition}

\begin{remark}[迹=测量=实际从属]\label{rmk:trace-measurement}
迹$\mathrm{Tr}(f)$把交互过程$F_2$变成非交互函数$F_1$——藏起自我更新$X\to X$，只留刺激-反应$A\to B$。这与量子测量（量子态$\to$经典数据）和马克思实际从属（活劳动$\to$死劳动/数据）是同一个数学操作。
\end{remark}

\begin{definition}[执行公式]\label{def:execution}
将$f:A\tensor X\to B\tensor X$写成分块算子（GoI I，Hilbert空间模型）：
$$f=\begin{pmatrix} f_{00}:A\to B & f_{01}:X\to B\\ f_{10}:A\to X & f_{11}:X\to X\end{pmatrix},$$
执行公式为
$$\mathrm{Ex}(f)=f_{00}+\sum_{n\geq0}f_{01}(f_{11})^n f_{10}.$$
$f_{11}:X\to X$是自我更新算子；$\sum_{n\geq0}(f_{11})^n$是反馈回路的迭代。
\end{definition}

\begin{definition}[余归纳执行]\label{def:coinductive-execution}
标准GoI要求$f_{11}$幂零（$(f_{11})^n=0$，反馈有限步终止），对应$\mu F_2$（有限交互）。对$\nu F_2$（生命），$f_{11}$不幂零——反馈永不终止但每步产生输出。\textbf{余归纳GoI}以productivity（每步有限时间内产出$b_n$）替换nilpotency作为正确性条件：
$$\mathrm{Ex}^\omega(f)=(b_0,b_1,b_2,\ldots)\in B^\omega,\qquad b_n=f_{01}(f_{11})^n f_{10}(a_n).$$
\end{definition}

\begin{theorem}[GoI语义证明：生命不可复制]\label{thm:goi-no-clone}\lvP
在GoI中，$\bang\nu F_2\to\nu\bang F_2$的解释需要复制反馈线$X$（收缩$X\to X\tensor X$），但$X$是线性状态（公理\ref{ax:no-contraction}），无收缩。因此不存在满足GoI连通性条件的算子解释。
\end{theorem}

\begin{proof}
$\nu F_2$的GoI网络中，反馈线$X$承载当前活状态——每步被$f_{11}$消耗并重建。$\bang\nu F_2$要求收缩整个网络，包括反馈线。但收缩$X\to X\tensor X$被公理\ref{ax:no-contraction}禁止。两个副本不能共享同一反馈线（那是同一过程而非两个副本），也不能分叉反馈线（$X$线性）。故不存在GoI解释。
\end{proof}

\begin{remark}[量论的GoI操作化路径]\label{rmk:measure-goi}
定义\ref{def:execution}给出量论参数的操作化路径（猜想，待在具体模型中验证）：
\begin{itemize}
  \item $T$（稳态/寿命）：$f_{11}$保持生产性的迭代次数（$\mu$情形为幂零指数，$\nu$情形为$\infty$或退化步数）；
  \item $N$（净方向）：$f_{11}$的谱半径$r(f_{11})$——$r<1$衰退，$r=1$稳态，$r>1$成长；
  \item $\alpha$（f-层级）：反馈嵌套深度（$f^1$无内部反馈，$f^2$一层，$f^3$含$X\to(X\to X)$高阶反馈）；
  \item $M=\alpha\cdot T\cdot N$：总交互量，即执行公式在寿命上的产出总和。
\end{itemize}
\end{remark}

\section{余归纳GoI：无限生产性过程}

标准GoI要求幂零性（$(f_{11})^n=0$，反馈有限步终止），对应$\mu F_2$。生命$\nu F_2$的反馈永不终止但每步产出。我们在$\mathbf{Rel}$中构造余归纳GoI。

\begin{definition}[流处理器]\label{def:stream-processor}
从$A$到$B$的流处理器是三元组$(X,R,x_0)$，其中$R\subseteq(A\times X)\times(B\times X)$是转移关系，$x_0\in X$是初始状态。它诱导因果关系$[\![ R]\!]\subseteq A^\omega\times B^\omega$：$(a_\bullet,b_\bullet)\in[\![ R]\!]$当且仅当存在$x_\bullet\in X^\omega$使对所有$n$，$((a_n,x_n),(b_n,x_{n+1}))\in R$。
\end{definition}

\begin{definition}[生产性]\label{def:productivity}
$R$是生产性的，如果对所有$(a,x)\in A\times X$，存在$(b,x')\in B\times X$使$((a,x),(b,x'))\in R$。
\end{definition}

\begin{theorem}[生产性保证完全性]\label{thm:productivity-totality}\lvP
若$R$生产性，则对每个$a_\bullet\in A^\omega$存在$b_\bullet\in B^\omega$使$(a_\bullet,b_\bullet)\in[\![ R]\!]$。
\end{theorem}

\begin{proof}
余归纳构造：由生产性，给定$(a_0,x_0)$存在$(b_0,x_1)$；再由生产性，给定$(a_1,x_1)$存在$(b_1,x_2)$；如此继续，构造完整的$b_\bullet$和$x_\bullet$。
\end{proof}

\begin{definition}[余归纳迹]\label{def:coinductive-trace}
余归纳迹$\mathrm{Tr}^\omega_X(R)=[\![ R]\!]\subseteq A^\omega\times B^\omega$将Mealy机映射为其无限输入输出行为。它满足迹公理的余归纳版本：
\begin{itemize}
  \item \textbf{vanishing}：$\mathrm{Tr}^\omega_I(R)$逐点应用；$\mathrm{Tr}^\omega_{X\times Y}=\mathrm{Tr}^\omega_X\circ\mathrm{Tr}^\omega_Y$
  \item \textbf{yanking}：$\mathrm{Tr}^\omega_X(\sigma)$诱导恒等流
  \item \textbf{superposing}：$g\otimes\mathrm{Tr}^\omega(f)=\mathrm{Tr}^\omega(g\otimes f)$
  \item \textbf{tightening}：$\mathrm{Tr}^\omega((g\otimes\mathrm{id})\circ f\circ(h\otimes\mathrm{id}))=g^\omega\circ\mathrm{Tr}^\omega(f)\circ h^\omega$
\end{itemize}
\end{definition}

\begin{theorem}[!不提升到流处理器]\label{thm:no-bang-stream}\lvP
不存在自然提升将生产性Mealy机$R\subseteq(A\times X)\times(B\times X)$映射为$\bang R\subseteq(\bang A\times\bang X)\times(\bang B\times\bang X)$使得$\bang R$的余归纳迹与$R$相容。
\end{theorem}

\begin{proof}
$\bang X$是$X$上的有限多重集。余归纳迹要求单一状态线索：每步消耗$x_n$产生$x_{n+1}$。多重集包含多个$X$元素，但元素之间没有身份（identity），无法确定哪个元素是"正在运行的当前状态"。$\bang R$的转移将多重集$\alpha$映射到$\beta$，但无法从$\beta$中追踪一条贯穿的状态线索。非交互$F_1(X)=X\tensor B$中$X$直接出现在输出中（不经过$\lolli$的codomain），配对唯一确定线索，故可提升（命题\ref{prop:lift-noninteractive}）；交互$F_2$中$X$是函数产出，多重集无法追踪。
\end{proof}

\begin{remark}[模型限定与一般化]\lvI
定理\ref{thm:no-bang-stream}的证明给出$\mathbf{Rel}$（关系范畴）中"多重集无法追踪单一状态线索"的构造性论证。此处"提升"指保持余归纳迹的函子映射$(R,x_0)\mapsto(\bang R,\bang x_0)$；其相对自然性以$\mathbf{Rel}$为底范畴。一般PTC中的完全一般性已由定理\ref{thm:bang-no-ptc}（猜想\ref{conj:bang-ptc}已解决）给出：$\kappa$-分叉、$\varepsilon$不唯一、$\bang$与$\nu$不交换三个论证不依赖任何特定模型。

本定理是种子定理\ref{thm:no-A-to-bangA}的GoI推广：种子定理说静态态射$A\to\bang A$不存在，本定理说动态流处理器的自然提升不存在。
\end{remark}

\begin{theorem}[执行不是沉积]\label{thm:execution-not-deposition}\lvP
GoI执行$\mathrm{Ex}(U)$或$\mathrm{Ex}^\omega(U)$不是$\mathcal{C}$中的态射，也不是$\bang$-模态对象。
\end{theorem}

\begin{proof}
有限情形$\mathrm{Ex}(U)=U_{00}+\sum_{n\geq0}U_{01}(U_{11})^nU_{10}$的结果是$\mathcal{C}$中态射，但计算这个和的过程是元层面操作（践演元预设（remark \ref{ax:performative}））。无限情形$\mathrm{Ex}^\omega(U)$产生$B^\omega$中的流，存在于流处理器范畴而非$\mathcal{C}$中。执行消耗资源（时间、活操作），而$\mathcal{C}$中态射不消耗任何东西——它只是存在。践演判断$\triangleright A$对应执行事件，不是类型。
\end{proof}

\begin{remark}[践演=执行]
矢列$\Gamma\vdash A$的GoI解释是静态连线$U$（阴，可被$\bang$复制存储）；践演判断$\triangleright A$对应$U$的执行$\mathrm{Ex}(U)$（阳，一次性事件）。你可以复制《资本论》（$\bang$），但不能复制"一个人读了《资本论》后改变生命"这个事件（$\triangleright$）。
\end{remark}

\section{生产性迹范畴：公理化}

将余归纳GoI从$\mathbf{Rel}$推广为一般范畴论概念。

\begin{definition}[生产性迹]\label{def:ptc}
设$\mathcal{C}$是SMC，对每个对象$A$有流对象$A^\omega=\nu X.(A\tensor X)$（最终余代数）。对Mealy机$(X,f:A\tensor X\to B\tensor X,x_0:I\to X)$，定义行为态射
$$\mathrm{beh}(f,x_0):A^\omega\to B^\omega$$
为余代数$c:X\tensor A^\omega\to B\tensor(X\tensor A^\omega)$的anamorphism（$c$在每步消耗输入头和当前状态，产生输出、新状态和剩余输入流）。$\mathcal{C}$是\textbf{生产性迹范畴}（PTC），如果$\mathrm{beh}$满足：
\begin{itemize}
  \item \textbf{Yanking}：$\mathrm{beh}(\sigma_{X,X},x_0)=\mathrm{id}_{X^\omega}$
  \item \textbf{Tightening}：$\mathrm{beh}((h\tensor\mathrm{id})f(g\tensor\mathrm{id}),x_0)=h^\omega\,\mathrm{beh}(f,x_0)\,g^\omega$
  \item \textbf{组合}：$\mathrm{beh}(f;g,(x_0,y_0))=\mathrm{beh}(g,y_0)\circ\mathrm{beh}(f,x_0)$
  \item \textbf{Productivity}：$f$完全则$\mathrm{beh}(f,x_0)$是完全因果流函数
\end{itemize}
\end{definition}

\begin{theorem}[$\mathbf{Rel}$是PTC]\label{thm:rel-ptc}\lvP
$\mathbf{Rel}$（笛卡尔积为张量）是生产性迹范畴。
\end{theorem}

\begin{proof}
以Tightening为例：$\mathrm{beh}((h\tensor\mathrm{id})f(g\tensor\mathrm{id}),x_0)$在$\mathbf{Rel}$中是关系复合$(h\times\mathrm{id})\circ R\circ(g\times\mathrm{id})$的余归纳闭包，等于$h^\omega\circ\mathrm{beh}(R,x_0)\circ g^\omega$（关系复合的单调性保证）。Yanking：$\sigma_{X,X}$在$\mathbf{Rel}$中是配对关系$((x_1,x_2),(x_2,x_1))$，其行为态射逐点交换流的头尾，余归纳地等于恒等流。Productivity：$\mathbf{Rel}$中完全关系（每输入至少有一个输出）保证每步有产出。
\end{proof}

\begin{theorem}[!不穿透生产性迹]\label{thm:bang-no-ptc}\lvP
在任意操作范畴中，对$F_2$型（交互型）Mealy机，不存在保持余归纳迹的自然$\bang$-提升。
\end{theorem}

\begin{proof}
\textbf{余归纳迹要求单一状态线索。}GoI分块中，Mealy机$f:A\tensor X\to B\tensor X$写为
$$f=\begin{pmatrix}f_{00}:A\to B & f_{01}:X\to B\\ f_{10}:A\to X & f_{11}:X\to X\end{pmatrix},$$
余归纳执行$\mathrm{Ex}^\omega(f)=(b_0,b_1,\ldots)$中，$f_{11}:X\to X$是\emph{单一状态转移}：每步$x_n$被消耗、$x_{n+1}=f_{11}(x_n)$产生、$x_{n+1}$成为下一步唯一当前状态。这条线索$x_0\to x_1\to x_2\to\cdots$是唯一的。

\textbf{$F_1$型的!-提升良定义。}对$F_1(X)=X\tensor B$（非交互），$f_{11}:X\to X$是直接状态传递（$X$在输出中直接出现，不经$\lolli$产出）。!-提升后$f^\bang_{11}=\bang f_{11}:\bang X\to\bang X$由余单子函子性给出，执行公式良定义（定理\ref{thm:nu-lift-data}）。

\textbf{$F_2$型的!-提升破坏单一线索。}对$F_2(X)=A\lolli(X\tensor B)$（交互），$f_{11}$不是直接状态传递：$X$在$\lolli$的后件（codomain）中，是函数的产出。!-提升后$f^\bang_{11}:\bang X\to\bang X$通过函数产出$\bang X$，但：

(1) $\bang X$享有收缩$\kappa:\bang X\to\bang X\tensor\bang X$（公理\ref{ax:bang}(ii)）。在反馈回路中，$\kappa$允许在任意步骤将状态分裂为两个副本：
$$\bang X\xrightarrow{f^\bang_{11}}\bang X\xrightarrow{\kappa}\bang X\tensor\bang X.$$
两个副本各自可独立继续反馈，产生两条并行的输出线索。余归纳迹要求\emph{单一}输出流$B^\omega$，但$\kappa$-分裂后产生$(B\tensor B)^\omega$——两条流而非一条。

(2) 要从$\bang X$中提取单一状态传回反馈，需要counit $\eps:\bang X\to X$（弃置/提取一次运行）。但$\eps$不保持线索完整性：$\eps\circ\kappa:\bang X\to X$选择一个副本丢弃另一个，且在$\mathbf{Rel}$中$\eps$是关系"多重集中的某个元素"而非函数——选择不唯一。在$\mathbf{Coh}$中$\eps$是clique包含，不提供唯一选择；在$\mathbf{Hilb}$中$\eps$要求严格压缩条件（$\|U_{22}\|<1$，定理\ref{thm:hilb-ptc}）。

(3) $\bang$与$\nu$不动点不交换：$\bang(A^\omega)\not\cong(\bang A)^\omega$。前者是"有限个同步无限流的沉积"（每步所有流同步演化），后者是"每步独立沉积的无限序列"（各步多重集独立变化）。定理\ref{thm:nu-lift-data}给出$\bang\nu F_1\to\nu(F_1\bang)$方向（$F_1$非交互），但逆方向$\nu(F_1\bang)\to\bang\nu F_1$不存在——一个每步独立变化的无限多重集序列不能被有限个同步流表示。对$F_2$，两个方向都不存在（定理\ref{thm:nu-no-lift-life}）。

\textbf{结论。}$F_2$型Mealy机的状态$X$在$\lolli$后件中被函数生产，!-提升后产出的$\bang X$是沉积（可复制、可丢弃、无规范当前元素）。余归纳迹要求单一状态线索，而$\kappa$允许分叉、$\eps$不提供唯一选择、$\bang\not\cong\nu\bang\nu$使类型不匹配。因此不存在保持余归纳迹的自然!-提升。\qed
\end{proof}

\begin{theorem}[有限与无限]\label{thm:finite-infinite}\lvP
若$\mathcal{C}$同时是TSMC和PTC：(a)幂零Mealy机（$(f_{11})^n=0$）的标准迹存在，行为在$n$步后退化为常数流；(b)生产性但非幂零的Mealy机标准迹不存在，但$\mathrm{beh}$存在。(a)对应$\mu F_2$（有限交互），(b)对应$\nu F_2$（生命）。
\end{theorem}

\begin{remark}[践演判断的范畴论地位]
$\mathrm{beh}(f,x_0):A^\omega\to B^\omega$不在$\mathcal{C}$中（流对象在余归纳GoI范畴$\mathcal{G}^\omega(\mathcal{C})$中）。构造$\mathrm{beh}$需要anamorphism，是范畴之间的操作而非范畴内部操作——这就是践演判断$\triangleright A$是元公理（践演元预设（remark \ref{ax:performative}））而非内部公理的原因：践演是你对形式系统\emph{做}的事，不是系统内部的命题。
\end{remark}

\section{资本与病毒：!-寄生体}

马克思区分两种流通：$C$-$M$-$C$（为买而卖，使用价值为目的，有限）和$M$-$C$-$M'$（为卖而买，交换价值为目的，无限），以及$M$-$M'$（生息资本，$G\to G+g$，最拜物教的形式）。

\begin{definition}[资本余代数]
资本对象$G$配备!-余代数$c:G\to\bang G$，通过contraction和dereliction得到复制$G\to\bang G\to\bang G\tensor\bang G\to G\tensor G$。
\end{definition}

\begin{theorem}[复制不是增长]\label{thm:copy-not-growth}\lvP
$G\to G\tensor G$把$g$映射为$(g,g)$，两个副本相同。这是复制（replication），不是增长（augmentation）：contraction保持值不变，只增加副本数量。
\end{theorem}

\begin{proof}
在$\mathbf{Rel}$中contraction是多重集对角嵌入$\alpha\mapsto(\alpha,\alpha)$，两分量完全相同。dereliction不增加任何东西。复合结果$g\mapsto(g,g)$，无增量。
\end{proof}

\begin{theorem}[增长必须消耗线性资源]\label{thm:growth-needs-linear}\lvP
在满足公理3(iii)（Seely同构$\bang(A\&B)\cong\bang A\tensor\bang B$）的操作范畴中，任何域完全由$\bang$-模态对象组成的态射只能产生$\bang$-模态对象（沉积），不能产生线性对象（不可复制的活资源）。形式化：不存在态射
$$\bang G\longrightarrow A$$
使$A$线性（无自然收缩$A\to A\tensor A$）且$A$持续更新。增长需要$A\tensor\bang G\to\bang G'\tensor B$型态射，其中线性资源$A$被消耗。
\end{theorem}

\begin{proof}
\textbf{第一步：Seely同构强制对角comultiplication。}余乘法$\delta_G:\bang G\to\bang G\tensor\bang G$通过Seely同构从笛卡尔对角$\Delta:G\to G\&G$（$g\mapsto(g,g)$）得到：
$$\bang G\xrightarrow{\bang\Delta}\bang(G\&G)\xrightarrow{\cong}\bang G\tensor\bang G.$$
因此$\delta_G$把一个沉积复制为两个\emph{相同}副本（由余交换律，两副本通过对称同构等价）。复制不产生新结构——两个副本携带完全相同的信息。

\textbf{第二步：!-模态态射只能复制/重排/丢弃。}在!-余代数的Kleisli范畴中，域为$\bang G$的态射由余代数结构（$\delta$复制、$\varepsilon$丢弃）和函子性（保持复合与恒等）生成。这些操作都不产生线性对象：$\delta$产出$\bang G\tensor\bang G$（仍是!-模态），$\varepsilon:\bang G\to G$产出$G$但$\bang G$被消耗（counit是余代数的"一次使用"，不提供$G$的持续更新），函子性只重排已有结构。

\textbf{第三步：线性对象的持续更新需要线性消耗。}若要产出持续更新的线性流$A^\omega=\nu X.(A\tensor X)$（每步有新的$A$），由定理\ref{thm:no-bang-stream}（!不穿透生产性状态），不存在$\bang G\to A^\omega$型态射——!-模态不能驱动无限线性流。每一步的$A$必须从外部线性资源获取并消耗：$A\tensor\bang G\to\bang G'\tensor B$。

\textbf{与定理\ref{thm:no-A-to-bangA}的对偶。}定理\ref{thm:no-A-to-bangA}说线性$A$不能被沉积为$\bang A$（生命不可资本化）；本定理说$\bang G$不能产生线性$A$（资本不能产生活劳动）。两个方向合在一起：线性资源与!-沉积之间没有自然的双向箭头，它们是根本不同的类型。\qed
\end{proof}

\begin{remark}[与马克思的对应]\lvI
本定理是马克思"活劳动是价值唯一源泉"的范畴论形式：$\bang G$（资本/死劳动）不能自我增殖，必须消耗$A$（活劳动/线性资源）才能增长。$M$-$C$-$M'$中的增量$M'-M$不来自$M$本身（$\bang G$只能复制），来自$C$中包含的活劳动$A$。"价值""剩余价值"的命名为[诠释]；数学内核是类型论的：!-模态不能产生线性资源。
\end{remark}

\begin{theorem}[!-模态增殖不是$F_2$型交互]\label{thm:fake-f2}\lvP
$G\multimap(G\tensor G')$型态射中，若$G$是$\bang$-模态（有余代数结构，不被消耗），则该态射不是$F_2$型交互。真正的$F_2$要求$X$线性（无收缩$X\to X\tensor X$，每步旧自我消耗、新自我产生）。$\bang$-模态的$G$不被消耗，其增殖是$G\to G+g$型（$\bang$-contraction），不是$S=f(S)$型（自我更新）。
\end{theorem}

\begin{remark}[{M-C-M\textquoteright 诠释}]\lvI
上述结构在政治经济学中对应马克思的$M$-$C$-$M'$（货币-商品-货币'）资本流通公式：投入的货币不被消耗，增殖是$\bang$-contraction。"资本"的命名为[诠释]。
\end{remark}

\begin{theorem}[两类!-寄生结构的范畴同构]\label{thm:virus-capital}\lvP
以下结构同构：!-余代数$P$+ 劫持自然变换$F_2\Rightarrow F_1$（定理\ref{thm:subsumption}）+ 沉积$P\tensor\bang B\to P'$。没有$F_2$宿主，$P$只能复制不能增长（定理\ref{thm:growth-needs-linear}）。
\end{theorem}

\begin{remark}[病毒与资本——诠释]\lvI
上述同构结构在生命论中被解释为病毒复制与资本增殖的同一。"病毒""资本"的命名为[诠释]。
\end{remark}

\begin{theorem}[!-寄生体随宿主削弱而增长枯竭]\label{thm:capital-contradiction}\lvP
$\bang$-寄生体须保存$F_2$才能增长（定理\ref{thm:growth-needs-linear}），但其趋势是把$F_2$转化为$F_1$（定理\ref{thm:subsumption}）；定理\ref{thm:no-bang-stream}保证$F_2$不能被完全转化。越接近完全转化，线性资源$A$的供给越少，增长越受限；极限$S=I$时$A=0$，无增长，寄生体自毁（定理\ref{thm:parasite-self-destruct}）。
\end{theorem}

\begin{proof}
增长需要$A$（来自$F_2$线性资源）；$F_2\to F_1$改造减少$A$供给；定理\ref{thm:no-bang-stream}保证$F_2$不被完全消灭，但越逼近极限，可提取的$A$越少；$S=I$时$A=0$，无增长，$\bang$-寄生体无宿主可寄生。\qed
\end{proof}

\begin{remark}[利润率趋向下降——诠释]\lvI
马克思"利润率趋向下降"被解释为本定理的政治经济学推论。"利润率"须在具体模型（如加权$\mathbf{Rel}$中的价值核算）中另行定义；此等同不作数学断言。
\end{remark}

\begin{remark}[$C$-$M$-$C$是生命过程]
$C$-$M$-$C$的结构是$A\tensor X\to B\tensor X'$：$A$（卖出的商品，线性，被让渡），$X$（卖者状态，线性，被消耗），$B$（买回的使用价值，线性，被消费），$X'$（需要满足后的新自我）。这是真正的$F_2$：有交互、有自我更新、过程有限（满足即止）。与$M$-$C$-$M'$的根本区别不在货币，在$X$是线性的（人被过程改变）而$G$是!-模态的（资本不被消耗）。
\end{remark}

\section{f-层级：明性的数学基础}

生命论区分三个递归层级：$f^1$（自在，自指维持）、$f^2$（自为，自我建模）、$f^3$（自觉，明性）。在GoI中，$\alpha$（量论$M=\alpha\cdot T\cdot N$中的f-层级）是反馈嵌套深度。

在compact closed范畴$\mathcal{G}(\mathcal{C})$中，每个态射$f:A\to B$有名字$\ulcorner f\urcorner:I\to A^*\tensor B$——过程的沉积表示。名字是$\bang$-模态的（可复制、存储、传输），过程本身是线性的。

\begin{definition}[f-层级]
$f^1$：一阶Mealy机$A\tensor X\to B\tensor X$，一条反馈线，系统维持自己但不看自己。
$f^2$：过程可访问自身名字$\bang\ulcorner f\urcorner$——反馈中嵌套自我模型，系统看着自己运行。
$f^3$：过程能区分运行$f$和名字$\ulcorner f\urcorner$——三条嵌套回路，看见模型是模型。
\end{definition}

\begin{theorem}[自我模型是!-模态的]\label{thm:self-model-bang}\lvP
任何可存储、检查、交流的自我模型必然是$\bang$-模态的。
\end{theorem}

\begin{proof}
存储、检查、交流要求复制（同一模型多次读取）和丢弃（模型可删除）。复制和丢弃是$\bang$的contraction和weakening。线性自我模型读取即消耗，不能被反思。故可反思的自我模型必须是$\bang$-模态的。
\end{proof}

\begin{corollary}[意识形态的存在论基础——诠释]\lvI
因为自我模型是$\bang$-模态的（定理\ref{thm:self-model-bang}），它可以从外部被植入、复制、传输。意识形态、文化、身份政治的运作基础正在于此。此条为哲学断言，其数学载体是定理\ref{thm:self-model-bang}；"意识形态"等同不作数学断言。
\end{corollary}

\begin{theorem}[运行被沉积模型接管时状态停止更新]\label{thm:alienation-mistake}\lvP
设$\bang\ulcorner f\urcorner$为$f^2$层可访问的沉积自我模型（$\bang$-模态，定理\ref{thm:self-model-bang}），$X$为线性运行状态。当$\bang\ulcorner f\urcorner$被错认为$X$本身时，系统按沉积模型运行自己（$S=f_{\bang}(S)$）而非自己运行自己（$S=f(S)$）；此后操作不产生新的$X'$，只重复$\bang$-模型的规定（$\bang$的纯收缩，空转）。
\end{theorem}

\begin{proof}
$X$是线性的（每步消耗、不可复制，定理\ref{thm:no-bang-stream}）；$\bang\ulcorner f\urcorner$是沉积的、固定的、可外部植入的；将后者等同于前者，运行权即交给模型，操作不产生新的$X'$。\qed
\end{proof}

\begin{remark}[异化——诠释]\lvI
上述结构在生命论中被命名为"异化"。此命名不作数学断言。
\end{remark}

\begin{theorem}[维持运行与沉积模型的区分]\label{thm:mingxing}\lvP
$f^3$层维持$X$（线性运行）与$\bang\ulcorner f\urcorner$（沉积模型）之间区分的能力：不是没有模型（$f^2$必要），是不把模型等同于自己。
\end{theorem}

\begin{proof}
$f^3$的第三条反馈回路将自我模型对象化：系统看到模型是$\bang$-模态的（固定的、可复制的、可能外部植入的），运行是线性的（每步更新、不可复制、自己的）。这个"看见"使系统能质疑被植入的模型、更新过时的模型、在模型和活操作间保持张力。\qed
\end{proof}

\begin{remark}[明性——诠释]\lvI
上述结构在生命论中被命名为"明性"。此命名不作数学断言。
\end{remark}

\begin{theorem}[$f^2$的哥德尔困境]\label{thm:godel-f2}\lvP
设$f^2$的转移函数$f_2:I_2\tensor S_2\to O_2\tensor S_2\tensor\bang\ulcorner f_2\urcorner$是图灵可计算的，且状态空间$S_2$可递归编码为自然数（$\ulcorner\cdot\urcorner:S_2\to\mathbb{N}$）。则：
\begin{enumerate}
  \item $f_2$的行为可在皮亚诺算术$\mathrm{PA}$中表示：存在公式$\mathrm{Trans}_2(x,y)$使得$\mathrm{PA}\vdash\mathrm{Trans}_2(\ulcorner s\urcorner,\ulcorner s'\urcorner)$当且仅当$f_2$把状态$s$转移到$s'$。
  \item 沉积的自我名字$\bang\ulcorner f_2\urcorner$在编码下恰好是$f_2$的哥德尔码$\ulcorner f_2\urcorner$；$\mathrm{PA}$可表示证明谓词$\mathrm{Proof}_T(x,y)$（"$x$是$y$在$T$中的证明的编码"）和可证性谓词$\mathrm{Prov}_T(y)=\exists x\,\mathrm{Proof}_T(x,y)$。
  \item 由对角线引理，存在哥德尔句子$G$使得$T\vdash G\leftrightarrow\neg\mathrm{Prov}_T(\ulcorner G\urcorner)$。若$T$一致则$T\nvdash G$且$T\nvdash\neg G$（G\"odel第一不完备）；若$T$一致则$T\nvdash\mathrm{Con}(T)$（第二不完备）。
  \item $f^3$的第三条反馈线（类型$\bang H_3\tensor X_3\to X_3$，$H_3=X_3\lolli X_3$）可将整个$T$对象化为$\bang\ulcorner T\urcorner\in\bang H_3$，从而在$T$\emph{之外}进行元层面推理。践演判断$\triangleright G$（活操作体知$G$为真）不是$T$内证明$\vdash G$，而是$f^3$运行中对$T$限度的直接看见——对应于$T$外的一致性证明（如Gentzen用$\varepsilon_0$超限归纳证明$\mathrm{PA}$一致，该证明不在$\mathrm{PA}$内）。
\end{enumerate}
\end{theorem}

\begin{proof}
(1) 图灵可计算函数在PA中可表示是数理逻辑基本结果（G\"odel 1931表示定理）：每个递归函数$f$存在公式$\phi_f$使$f(\vec{n})=m\Rightarrow\mathrm{PA}\vdash\phi_f(\vec{n},m)$且$f(\vec{n})\neq m\Rightarrow\mathrm{PA}\vdash\neg\phi_f(\vec{n},m)$。$f_2$图灵可计算故$\mathrm{Trans}_2$存在。

(2) $\bang\ulcorner f_2\urcorner$是$f_2$代码的沉积（!-模态），在自然数编码下就是哥德尔码$\ulcorner f_2\urcorner$。证明谓词$\mathrm{Proof}_T(x,y)$是原始递归的（检查$x$是否为$y$的合法证明序列），故在PA中可表示。

(3) 对角线引理：对任何只含一个自由变元的公式$\psi(x)$，存在句子$\phi$使$T\vdash\phi\leftrightarrow\psi(\ulcorner\phi\urcorner)$。取$\psi(x)=\neg\mathrm{Prov}_T(x)$得$G$。若$T\vdash G$则$T\vdash\mathrm{Prov}_T(\ulcorner G\urcorner)$（可证性谓词性质），又$T\vdash G\leftrightarrow\neg\mathrm{Prov}_T(\ulcorner G\urcorner)$故$T\vdash\neg\mathrm{Prov}_T(\ulcorner G\urcorner)$，矛盾（若$T$一致）。若$T\vdash\neg G$则$T\vdash\mathrm{Prov}_T(\ulcorner G\urcorner)$，$\omega$-一致时$T\vdash G$，矛盾。第二不完备：$\mathrm{Con}(T)=\neg\mathrm{Prov}_T(\ulcorner 0=1\urcorner)$在$T$内不可证。

(4) $f^3$的自我函数空间$H_3=X_3\lolli X_3$包含所有自函数（定理\ref{thm:f3-fixedpoint}证明），包括"在$T$中推理"这个函数$h_T:X_3\to X_3$。$f^3$可读取$\bang h_T$（沉积的$T$-推理规则）并\emph{选择不按它运行}——这是$f^2$做不到的：$f^2$读取$\bang\ulcorner f\urcorner$后直接按模型运行，无"判断"步骤。$f^3$的判断步骤对应于元层面：它能把$T$本身作为对象$\bang\ulcorner T\urcorner$审视，从而看见$G$为真（$G$说"我在$T$中不可证"，而$f^3$在$T$外确实看见$T$证明不了$G$）。践演判断$\triangleright G$是这个"看见"的范畴论表达：不是$T$内的矢列$\vdash G$，是活操作对$T$限度的直接体知。Gentzen证明PA一致用了$\varepsilon_0$超限归纳——这个推理不在PA内（否则PA证明自己一致，违反第二不完备），但它是一个有效的数学操作；同理，$\triangleright G$不在$T$内但在$f^3$运行中有效。\qed
\end{proof}

\begin{remark}[模型条件与新颖性界分]\lvI
本定理的数学内核（(1)-(3)）是G\"odel不完备定理的直接应用，不是新结果——本文不掠前人之美。新颖之处仅在(4)的范畴论对应：$f^2/f^3$区分精确对应于"系统内/系统外"区分，$\bang\ulcorner f\urcorner$对应哥德尔码，践演判断$\triangleright A$对应系统外的元层面操作。这把生命论的f-层级理论与证明论中"形式系统不能自证其一致性"的经典结果连接起来。\end{remark}

\begin{remark}[$\alpha$的GoI定义与记号对应]
在GoI执行公式中，$\alpha$是反馈算子$f_{11}$（GoI分块记号，对应PTC/量论中的$U_{22}$）的嵌套深度：$f^1$有一层迭代$\sum(f_{11})^n$，$f^2$的$f_{11}$本身含反馈（迭代中的迭代），$f^3$再深一层。

\textbf{两套分块记号对应}：GoI（定义\ref{def:execution}）用$f_{00}/f_{01}/f_{10}/f_{11}$，PTC/量论（定义\ref{def:ptc}、§20）用$U_{11}/U_{12}/U_{21}/U_{22}$；对应关系为$f_{00}\leftrightarrow U_{11}$、$f_{01}\leftrightarrow U_{12}$、$f_{10}\leftrightarrow U_{21}$、$f_{11}\leftrightarrow U_{22}$（反馈算子）。
\end{remark}

\section{可靠性与量子对应}

\begin{theorem}[可靠性]\label{thm:soundness}\lvP
如果$\pi$是余归纳片段$\mu LL^\nu$中$\Gamma\vdash A$的证明，则其GoI解释$U(\pi)$是生产性流处理器。
\end{theorem}

\begin{proof}[证明（概述；完整结构归纳见补充材料）]
对证明结构归纳。公理解释为恒等连线（生产性）；张量规则为并列/拆分（保持生产性）；$\multimap L$是Mealy机复合（定理\ref{thm:rel-ptc}中复合保持生产性）；Cut是反馈连接（同复合）；$\nu R$（余迭代）把一步转移展开为无限流（余归纳构造保证每步有产出）；$\nu L$取流的头（生产性流保证头存在）；!规则只操作$\bang$-模态数据（不涉及线性状态消耗）。每种规则保持生产性。
\end{proof}

\begin{corollary}[诠释]\lvI
操作$\nu F_2$型资源的证明不能被promotion（不能变成可复制的教条）——这是定理\ref{thm:no-bang-stream}的证明论对应：关于生命过程的证明本身必须是"活的"，不能冻结为可复制的定理。"生命过程的证明"的等同为[诠释]，其数学载体（$\nu F_2$不可提升）为[定理]。

若需保留纯证明论陈述，可另写一条[定理]：在$\mu$MALL+$\mu$中，promotion规则不可作用于以（自由）线性原子为主公式的证明（证明见定理\ref{thm:no-bang-stream}/\ref{thm:nu-no-lift-life}）。
\end{corollary}

\begin{theorem}[不可克隆是推论]\label{thm:no-cloning-corollary}\lvP
在$\mathbf{FdHilb}$（有限维Hilbert空间与线性映射）中，量子对象（维数$\geq2$）无自然收缩$\Delta:A\to A\tensor A$，量子过程构成生产性迹范畴。Hilb$^\omega$满足定理\ref{thm:bang-no-ptc}的状态唯一可追踪性条件（量子态有规范的正交分解，测量给出唯一经典输出，见§21与定理\ref{thm:hilb-ptc-full}）。由定理\ref{thm:bang-no-ptc}，$\bang$不提升到量子交互流处理器——这就是Wootters-Zurek不可克隆定理。
\end{theorem}

\begin{proof}
量子态是线性的（不可克隆，Wootters-Zurek 1982）；幺正演化是生产性Mealy机（每步产生新状态）；量子测量是迹（藏起量子态，产生经典数据）。$\mathbf{FdHilb}$以因果流函数为行为范畴构成PTC。定理\ref{thm:bang-no-ptc}直接给出不可克隆。
\end{proof}

\begin{remark}[测量=迹=实际从属]
量子测量（量子态$\to$经典数据）、GoI迹（藏起反馈线$X$）、马克思实际从属（活劳动$\to$死劳动）是同一个数学操作：藏起线性状态线索，只留可复制的输入输出行为。三者同一不是比喻，是范畴论结构的同一。
\end{remark}

\begin{theorem}[外层$\bang$不穿透的$\mu$情形严格化]\label{thm:mu-rigorous}\lvP
$\bang\mu F_2\to\mu\bang F_2$在IMELL+$\mu$中不可证。穷尽$\otimes R$的所有拆分：非循环分支需要从空上下文证明线性资源（不可能）或promotion线性原子（不可能）；循环分支中$\mu$-展开不降低矢列复杂度（$\mu F_2$展开为$A\multimap(\mu F_2\tensor B)$，$\mu F_2$重新出现），证明树无限生长但不产生公理叶子。
\end{theorem}

\section{$G^\omega(\mathcal{C})$：生产性紧闭合范畴}

\begin{definition}[$G^\omega(\mathcal{C})$]\label{def:gomega}
设$\mathcal{C}$是PTC。$G^\omega(\mathcal{C})$的对象是$(A^+,A^-)$（$\mathcal{C}$中对象的有序对），态射$(A^+,A^-)\to(B^+,B^-)$是$\mathcal{C}$中的生产性Mealy机$A^+\tensor B^-\tensor X\to B^+\tensor A^-\tensor X$，复合用生产性迹$\mathrm{Tr}^\omega$。
\end{definition}

\begin{theorem}[$G^\omega(\mathcal{C})$是dagger紧闭合范畴]\label{thm:gomega-compact}\lvP
$G^\omega(\mathcal{C})$是范畴，且是dagger compact closed范畴：对偶$(A^+,A^-)^*=(A^-,A^+)$，unit/counit由$\mathcal{C}$中的弯线给出，蛇方程由Yanking公理保证（类比Joyal--Street--Verity 1996的Int-构造）。$G^\omega(\mathcal{C})$还有递归类型：流对象$A^\omega=\nu X.(A\tensor X)$是合法对象（公理4），$\nu F_2$（生命流类型）在$G^\omega(\mathcal{C})$中可解释。
\end{theorem}

\begin{proof}
复合封闭由Productivity公理；结合律由Tightening公理；单位律由Yanking公理。紧闭合结构直接类比Int-构造（Joyal-Street-Verity 1996），但用生产性迹代替普通迹。递归类型来自PTC的流对象和余归纳迹。
\end{proof}

\begin{theorem}[!在$G^\omega(\mathcal{C})$中仍不穿透]\label{thm:bang-gomega}\lvP
在$G^\omega(\mathcal{C})$中不存在$\bang\nu F_2\to\nu\bang F_2$。$G^\omega(\mathcal{C})$中的态射是$\mathcal{C}$中的生产性Mealy机，若该提升存在则底层$\mathcal{C}$中也存在，与定理\ref{thm:bang-no-ptc}矛盾（定理\ref{thm:bang-no-ptc}在任意操作范畴中成立，无需额外模型条件）。
\end{theorem}

\begin{remark}[与种子定理同族]
本定理与种子定理\ref{thm:no-A-to-bangA}（不存在$A\to\bang A$）同族：$G^\omega$中的$\bang$穿透若存在，将诱导底层$\mathcal{C}$中的$\bang$-结构，与定理\ref{thm:bang-no-ptc}矛盾。
\end{remark}

\begin{remark}[与Int-构造的关系]
$\mathrm{Int}(\mathcal{C})$是$G^\omega(\mathcal{C})$的全子范畴，由无状态（$X=I$）态射组成。Int(C)处理有限连线（nilpotent反馈），$G^\omega(\mathcal{C})$处理无限流处理器（productive反馈）。$G^\omega$是Int构造的"无限版"。
\end{remark}

\begin{remark}[PTC实例]
Rel$^\omega$（关系与流）是已验证的PTC实例。Hilb$^\omega$（Hilbert空间与因果有界算子，生产性迹=Neumann级数$\sum U_{22}^n$，收敛条件$\|U_{22}\|<1$，即控制论小增益定理）和Coh$^\omega$（相干空间与因果clique）是合理的候选实例，严格验证待补。三个实例中!不穿透线性状态都成立——这是PTC公理的普遍推论，不是特定模型的偶然性质。
\end{remark}

\section{f-层级的封闭性}

\begin{theorem}[$f^3$是不动点]\label{thm:f3-fixedpoint}\lvP
在$G^\omega(\mathcal{C})$中，层级提升操作$L$（增加一条反馈嵌套）在第三层达到不动点：$L^3=L^n$对所有$n\geq3$。即$f^4$层可定义的任何操作，$f^3$层已可定义。
\end{theorem}

\begin{proof}
$f$-层级对应GoI反馈线的嵌套深度。设$X_n$为$f^n$层的状态类型，$H_n=X_n\lolli X_n$为自函数空间。

\textbf{前两层是质变。}$f^1$有一条反馈线（状态更新$X_1\to X_1$），无自我模型。$f^2$增加第二条反馈线，可访问自身沉积名字$\bang\ulcorner f\urcorner$（类型$\bang H_2$），系统能"看自己运行"。$f^3$增加第三条反馈线，类型为$\bang H_3\tensor X_3\to X_3$：读取一个沉积的自我函数$\bang h$和当前状态$x$，产生新状态$h(x)$。

\textbf{第三条反馈线已包含任意高阶变换。}关键观察：$H_3=X_3\lolli X_3$是$X_3$上\emph{所有}自函数的空间，包含：
\begin{enumerate}
  \item 任何具体运行策略$h:X_3\to X_3$；
  \item 任何模型变换$g:H_3\to H_3$：定义$\mathrm{enc}(g):X_3\to X_3$为
  $$\mathrm{enc}(g)(x)=\mathrm{let}\,(\bang h, x')=\pi(x)\,\mathrm{in}\,\iota(g(h), x'),$$
  其中$\pi:X_3\to\bang H_3\tensor X_3$从状态中提取当前沉积模型与剩余状态（由$X_3\cong\cdots\tensor\bang H_3\tensor\cdots$的out映射给出），$\iota$为其逆。$\mathrm{enc}(g)$是良定义的$X_3\to X_3$态射；
  \item 任意$k$阶变换$g:H_3^{(k)}\to H_3^{(k)}$（其中$H_3^{(0)}=X_3$，$H_3^{(k+1)}=H_3^{(k)}\to H_3^{(k)}$），通过重复$\enc$编码为$H_3$中的元素。
\end{enumerate}

\textbf{$f^4$不增加新能力。}$f^4$的第四条反馈线类型为$\bang(H_3\lolli H_3)\tensor X_4\to X_4$，即读取沉积的模型变换$\bang g$并应用。但由(2)，任何$g:H_3\to H_3$都被$\mathrm{enc}(g)\in H_3$编码，第三条反馈线读取$\bang\mathrm{enc}(g)$即可实现$g$的功能。编码$\enc$是单射：给定$\mathrm{enc}(g)$，取任意$h\in H_3$和$x_0\in X_3$，$g(h)=\pi_1(\mathrm{enc}(g)(\iota(h,x_0)))$可恢复$g$。因此第四条反馈线的信息可无损编码到第三条中。

由归纳，所有$n\geq3$的$f^n$层操作等价于$f^3$。$f^1\to f^2$（无模型→有模型）和$f^2\to f^3$（被模型支配→能对象化模型）是质变；$f^3\to f^4$不是质变，因为$f^3$的$H_3$已穷尽所有自指变换。\qed
\end{proof}

\begin{remark}[$\alpha$只取三个值]
$\alpha\in\{1,2,3\}$：自在、自为、自觉。不存在"比明性更明"的状态，只存在明性操作得更持续、更彻底。这对应禅宗三段（见山是山$\to$见山不是山$\to$见山还是山），第三段之后没有第四段。
\end{remark}

\section{量论操作化}

在加权关系模型中，$U_{22}:X\to X$是非负矩阵。量论参数获得精确计算定义：

\begin{definition}[量论参数]
~
\begin{itemize}
\item $T$（寿命）：有限过程$T=\min\{n:U_{22}^n=0\}$（幂零指数）；无限过程$T=\infty$。
\item $N$（净方向）：$N=\rho(U_{22})$，$U_{22}$的Perron-Frobenius谱半径。$N<1$异化，$N=1$稳态，$N>1$创造。
\item $\alpha$（f-层级）：反馈嵌套深度，$\alpha\in\{1,2,3\}$。
\item $M$（意义总量）：有限过程$M=\alpha\cdot T\cdot N$；无限过程$M=\alpha\sum_{n=0}^\infty N^n$。
\item $W=O/M$（萃取率）：被!-系统取走的产出占总意义的比例。
\item $l'=(1-N)\times100\%$（异化率，$N\leq1$时）。
\item $C=R-A$（革命意识）：认清$R$减去接受$A$。
\end{itemize}

\noindent\textbf{模型内定义。}在加权$\mathbf{Rel}$模型中，$M$有独立于哲学诠释的纯数学定义：设$U_{22}$为非负转移矩阵，$M$是状态转移图中所有路径的总权重（有限过程为$\alpha\sum_{n=0}^{T-1}\rho(U_{22})^n$，无限过程为$\alpha\sum_{n=0}^{\infty}\rho(U_{22})^n$）。这是一个可计算的纯数学量。将$M$命名为"意义总量"是[诠释]，不是定义的一部分。
\end{definition}

\begin{theorem}[几何级数收敛]\label{thm:slow-death}\lvP
当$N<1$时，即使过程永远运行（生产性），加权和$\alpha\sum_{n=0}^\infty N^n=\alpha/(1-N)<\infty$。
\end{theorem}

\begin{remark}[诠释："慢性死亡"]\lvI
在量论诠释下（$M$=意义总量，$N$=净方向），上述收敛意味着：$N<1$的系统虽然永远运转，但意义总量有限——每步产出按几何级数衰减。"慢性死亡"是这一数学事实的哲学命名，不是数学推导。
\end{remark}

\begin{theorem}[萃取加速异化]\label{thm:extraction-alienation}\lvP
萃取（在反馈回路中插入$\bang$-投影$\pi_!$，把线性状态沉积化）在两个层面上加速异化：
\begin{enumerate}
  \item \textbf{一般PTC层面}：由定理\ref{thm:bang-no-ptc}，$\bang$-投影破坏余归纳迹的单一输出流（$\kappa$-分叉论证）。萃取比例越高，反馈回路中被沉积化的线性状态越多，生产性迹越弱；完全萃取（$W=1$，所有线性状态被$\pi_!$捕获）时反馈回路中无线性状态，生产性迹不存在，过程终止。
  \item \textbf{加权$\mathbf{Rel}$量化}：萃取$W$对应反馈算子缩放$U_{22}\mapsto(1-W)U_{22}$，故$N(W)=\rho((1-W)U_{22})=(1-W)N(0)$。$W$越高$N$越低；$W=1$时$N=0$，过程立即死亡（$S\to\I$）。
\end{enumerate}
\end{theorem}

\begin{proof}
(1) 萃取操作在GoI反馈回路中插入$\pi_!:X\to\bang X$（把线性状态投影到沉积模型），复合后反馈算子变为$U_{22}\circ\pi_!$（或等价地，回路中插入$\bang$-收缩）。由定理\ref{thm:bang-no-ptc}的$\kappa$-分叉论证，$\bang$-收缩允许状态分裂为两个副本，破坏余归纳迹要求的单一输出流。萃取比例$W$对应$\pi_!$的作用范围：$W=0$时$\pi_!$为恒等（无沉积化），$W=1$时$\pi_!$捕获全部线性状态（无线性资源参与反馈）。$W=1$时反馈回路中无$X$（线性状态），$U_{22}=0$，迹不存在，过程终止。

(2) 在加权$\mathbf{Rel}$中，$\pi_!$实现为反馈缩放$(1-W)U_{22}$（比例$W$的状态被沉积化、不参与反馈），由谱半径齐次性$\rho(cU)=c\rho(U)$（$c\geq0$）得$N(W)=(1-W)N(0)$。\qed
\end{proof}

\begin{theorem}[明性抵抗磨损]\label{thm:mingxing-resists}\lvI
$C>0$（认清多于接受）时$N$趋向$\geq1$；$C<0$时$N$趋向$<1$。f³运行时线性过程$X$保持自我更新；f²把$\bang\ulcorner f\urcorner$当自己时$U_{22}$被$\bang$-收缩削弱。
\end{theorem}

\begin{theorem}[动态自毁]\label{thm:dynamic-self-destruct}\lvC
在加权$\mathbf{Rel}$模型中，设反自指寄生体每步萃取比例$\beta(t)\in[0,1]$的线性资源（反馈状态被沉积化的比例），则：
\begin{enumerate}
  \item $N(t)=\rho(U_{22}(t))=N(0)\prod_{s=0}^{t-1}(1-\beta(s))$单调不增；
  \item 若存在$\beta_0>0$使$\beta(t)\geq\beta_0$对所有$t$（持续萃取不衰减），则$N(t)\leq N(0)(1-\beta_0)^t\to0$（指数衰减）；
  \item 若$\beta(t)\to0$但$\sum_{t=0}^\infty\beta(t)=\infty$（萃取率递减但总量发散），则$N(t)\leq N(0)\exp(-\sum_{s=0}^{t-1}\beta(s))\to0$；
  \item $N(t)\to0$时$M(t)=\alpha\sum_{n=0}^\infty N(t)^n\to\alpha$（意义总量趋向f-层级系数本身，无增量），$S\to\I$，寄生体无宿主可寄生（定理\ref{thm:parasite-self-destruct}）。
\end{enumerate}
\end{theorem}

\begin{proof}
(1) 每步萃取$\beta(t)$比例后，反馈算子缩放为$U_{22}(t+1)=(1-\beta(t))U_{22}(t)$（加权$\mathbf{Rel}$中萃取=反馈缩放，定理\ref{thm:extraction-alienation}）。由谱半径齐次性$\rho(cU)=|c|\rho(U)$（$c\geq0$），$N(t+1)=(1-\beta(t))N(t)$，归纳得$N(t)=N(0)\prod_{s=0}^{t-1}(1-\beta(s))$。因$1-\beta(s)\leq1$，$N(t)$单调不增。

(2) $\beta(t)\geq\beta_0>0$时$N(t)\leq N(0)(1-\beta_0)^t$，由$0<1-\beta_0<1$得指数衰减至0。

(3) 由不等式$1-x\leq e^{-x}$（$x\geq0$），$\prod_{s=0}^{t-1}(1-\beta(s))\leq\exp(-\sum_{s=0}^{t-1}\beta(s))$。$\sum\beta(s)=\infty$时右端趋向0。

(4) $N(t)\to0$时，$M(t)=\alpha/(1-N(t))\to\alpha$（$N<1$时几何级数和），即意义总量只剩f-层级系数的常数——无成长、无盈余、无新意义。反馈算子$U_{22}(t)\to0$意味着无状态更新，$S\to\I$，由定理\ref{thm:parasite-self-destruct}寄生体退化为纯收缩（空转），无增长。\qed
\end{proof}

\begin{remark}[抵抗条件]\lvC
动态自毁是\emph{无抵抗条件下}的结论。若f³明性操作（$C>0$）每步恢复比例$\gamma(t)$的线性资源（把沉积状态重新激活为线性运行），则有效萃取率为$\beta(t)-\gamma(t)$。当$\gamma(t)\geq\beta(t)$时$N(t)$不再下降；当$\gamma(t)>\beta(t)$时$N(t)$回升。革命（定理\ref{thm:revolution-cut}）是$\gamma=1$的极端情形：一次性消除!-寄生Cut，$N$完全恢复。\end{remark}

\begin{remark}["资本主义必然灭亡"的数学形式]\lvI
本定理给出"反自指寄生体必然自毁"的动态形式：不是静态结构矛盾（定理\ref{thm:capital-contradiction}），而是\emph{持续萃取必然导致宿主枯竭}的指数衰减律。但"必然"不意味着"自动"——衰减可以很慢（$\beta_0$很小），且寄生体可以通过地理扩张、技术创新、新市场开发等方式暂时补充线性资源$A$（提高$N(0)$或重置$t$），延缓自毁。抵抗条件（上述remark）说明：自毁不是宿命，是操作——要么寄生体持续萃取直到宿主枯竭，要么活人通过f³明性和革命切断萃取回路。
\end{remark}

\section{完全性}

\begin{theorem}[有限片段GoI完全性——已知结果]\label{thm:completeness-finite}\lvP
在有限非交互片段（$F_1$型，$\mu$不动点）中，GoI解释对MELL（乘法指数线性逻辑）片段完全：每个nilpotent连线对应一个证明。此为Girard（1989）与Danos--Regnier（1995）的经典结果，本文直接引用，不重新证明。
\end{theorem}

\begin{theorem}[生命过程超越有限证明]\label{thm:life-beyond-proof}\lvP
设$A$、$B$、$X$为可数无限集合。在余归纳片段（$\nu F_2$型）中，存在生产性Mealy机不对应任何有限证明。
\end{theorem}

\begin{proof}
生产性Mealy机$(X,R,x_0)$由转移关系$R\subseteq(A\times X)\times(B\times X)$决定。当$A,B,X$可数时，$(A\times X)\times(B\times X)$可数，故$R$的集合基数为$2^{\aleph_0}$（连续统）。生产性条件（对所有$(a,x)$存在$(b,x')$）只排除部分关系，满足生产性的$R$的集合基数仍为$2^{\aleph_0}$（对每个$(a,x)$至少选一个$(b,x')$，选择函数的数量为$\aleph_0^{\aleph_0}=2^{\aleph_0}$）。

有限证明是有限符号串，符号表可数，故有限证明的集合基数为$\aleph_0$（可数）。

$2^{\aleph_0}>\aleph_0$（Cantor对角线论证），故存在生产性Mealy机不对应任何有限证明。\qed
\end{proof}

\begin{remark}[证明论强度的边界]
定理\ref{thm:life-beyond-proof}不依赖于任何特定形式系统的不完备性（如G\"odel不完备需要足够强的算术），只依赖于\emph{有限性}本身：有限证明是可数的，生产性过程是连续统。这是一个比G\"odel更基础的限制——不是系统不够强，是有限手段原则上不能穷尽无限生产性过程。
\end{remark}

\begin{remark}[践演比证明根本]
定理\ref{thm:life-beyond-proof}是"理论不能代替实践"的证明论表达：有限证明（理论）能保证结构性正确（可靠性），但不能穷尽无限生产性过程（践演）。有些真理只能通过"活出来"来知道。这是哥德尔不完备的范畴论版本，f³/践演是出路。
\end{remark}

\section{Hilb$^\omega$：量子PTC实例}

\begin{theorem}[Hilb$^\omega$是PTC]\label{thm:hilb-ptc}\lvP
以可分Hilbert空间为对象、因果有界线性算子为态射、$\ell^2$流对象、严格压缩反馈（Neumann级数$\sum U_{22}^n$，$\|U_{22}\|<1$时收敛）构成PTC。
\end{theorem}

\begin{proof}
SMC结构标准（Hilbert张量积）。流对象$\ell^2(\mathbb{N};A)$有等距cons和hd/tl。生产性迹$\mathrm{Tr}^\omega(U)=U_{11}+U_{12}(I-U_{22})^{-1}U_{21}$在$\|U_{22}\|<1$时有界（Neumann级数定理）。Yanking在无实际反馈（$U_{12}=0$或$U_{21}=0$）时为$U_{11}$。Tightening是矩阵分块乘法的推论。复合保持严格压缩：$\|U_{22}V_{22}\|\leq\|U_{22}\|\|V_{22}\|<1$。Productivity由Neumann级数收敛保证。
\end{proof}

\begin{theorem}[Hilb$^\omega$完整PTC]\label{thm:hilb-ptc-full}\lvP
以每步有界的因果流$A^{\bar\omega}$（$\ell^\infty$型）为流对象、幂有界反馈（$\sup_n\|U_{22}^n\|<\infty$）为生产性条件，Hilb$^\omega$是PTC。$N<1$时输出$\ell^2$有界（衰减/暂态）；$N=1$时输出每步有界（持续/稳态，覆盖酉演化）；$N>1$时输出增长（创造）。
\end{theorem}

\begin{proof}
幂有界保证每步输出$b_n=U_{11}(a_n)+\sum_{k<n}U_{12}U_{22}^{n-1-k}U_{21}(a_k)$有界。幂有界算子复合仍幂有界。其余公理同定理\ref{thm:hilb-ptc}。
\end{proof}

\begin{theorem}[Coh$^\omega$是PTC]\label{thm:coh-ptc}\lvP
以相干空间（Girard 1987）为对象、因果线性映射（clique）为态射、clique流为流对象，Coh$^\omega$是PTC。
\end{theorem}

\begin{proof}
\textbf{预备。}相干空间$\mathcal{A}$由web $|\mathcal{A}|$和自反对称相干关系$\coh_{\mathcal{A}}$组成；clique是$|\mathcal{A}|$的子集，其中任意两元素相干。$\mathbf{Coh}$的SMC结构是标准的：$\mathcal{A}\otimes\mathcal{B}$的web为$|\mathcal{A}|\times|\mathcal{B}|$，$(a,b)\coh(a',b')$当且仅当$a\coh a'$且$b\coh b'$；线性蕴含$\mathcal{A}\lolli\mathcal{B}$中$(a,b)\coh(a',b')$当且仅当$a\coh a'$蕴涵$b\coh b'$。

\textbf{流对象。}$\mathcal{A}^\omega=\nu X.(\mathcal{A}\otimes X)$在$\mathbf{Coh}$中存在（Girard不动点构造）：web为$|\mathcal{A}|$的有限非空序列$|\mathcal{A}|^+$，$s\coh t$当且仅当在公共长度上逐点相干。clique对应于无限流的有限前缀闭包（相干树）。

\textbf{Mealy机的clique条件。}生产性Mealy机$f$是$(\mathcal{A}\otimes X)\lolli(\mathcal{B}\otimes X)$中的clique，其元素为$(a,x,b,x')$，表示"在状态$x$接收$a$，产出$b$，转移到$x'$"。clique条件为：
$$(a,x)\coh(a',x')\implies(b,x')\coh(b',x'').$$
生产性（完全性）：对所有$(a,x)$存在$(b,x')$使$(a,x,b,x')\in f$。

\textbf{生产性迹保持clique。}设$\mathrm{beh}(f,x_0)$诱导的流关系为：$(a_\bullet,b_\bullet)\in\mathrm{beh}(f,x_0)$当且仅当存在$x_\bullet$使$(a_n,x_n,b_n,x_{n+1})\in f$对所有$n$。

需证$\mathrm{beh}(f,x_0)$是$\mathcal{A}^\omega\lolli\mathcal{B}^\omega$中的clique。设$(a_\bullet,b_\bullet),(a'_\bullet,b'_\bullet)\in\mathrm{beh}(f,x_0)$，分别由状态序列$x_\bullet,x'_\bullet$见证。假设$a_\bullet\coh a'_\bullet$（逐点相干）。由$x_0=x'_0$（同一初始状态，相干性自反），若$x_n\coh x'_n$，则$(a_n,x_n)\coh(a'_n,x'_n)$，由$f$的clique条件得$(b_n,x_{n+1})\coh(b'_n,x'_{n+1})$，即$b_n\coh b'_n$且$x_{n+1}\coh x'_{n+1}$。归纳得对所有$n$，$b_n\coh b'_n$，即$b_\bullet\coh b'_\bullet$。故$\mathrm{beh}(f,x_0)$是clique。

\textbf{PTC公理。}
\emph{Yanking}：$\sigma_{X,X}$的clique条件由$\coh$的对称性保证，$\mathrm{beh}(\sigma,x_0)$逐点交换，余归纳地等于恒等流clique。
\emph{Tightening}：前处理$g$和后处理$h$都是clique，clique在$\mathbf{Coh}$的复合下封闭，故$\mathrm{beh}((h\otimes\mathrm{id})f(g\otimes\mathrm{id}),x_0)=h^\omega\circ\mathrm{beh}(f,x_0)\circ g^\omega$。
\emph{组合}：$f$的输出clique与$g$的输入clique复合，clique复合封闭，故$\mathrm{beh}(f;g,(x_0,y_0))=\mathrm{beh}(g,y_0)\circ\mathrm{beh}(f,x_0)$。
\emph{Productivity}：$f$完全（生产性）保证每步至少有一个$(b,x')$，归纳构造给出完整的$b_\bullet$和$x_\bullet$；上述clique证明保证输出流相干。\qed
\end{proof}

\begin{remark}[三个PTC实例]
Rel$^\omega$（经典计算）、Hilb$^\omega$（量子力学）、Coh$^\omega$（证明论）三个实例中$\bang$不穿透生产性交互状态都成立。这从计算、物理、逻辑三个领域验证了同一个结构事实：线性状态的交互过程不可复制。
\end{remark}

\section{革命跃迁}

\begin{theorem}[异化降低$N$]\label{thm:alienation-lowers-n}\lvC
异化态的反馈算子为$U_{22}\circ\pi_!$，其中$\pi_!$是通过!-模型的投影（$\|\pi_!\|\leq1$，非满射时丢失线性状态的不可复制部分）。在加权$\mathbf{Rel}$模型中：
\begin{enumerate}
  \item $\rho(U_{22}\circ\pi_!)\leq\rho(U_{22})$恒成立；
  \item 若$U_{22}$不可约非负且$\pi_!\neq\mathrm{id}$，且Perron特征向量$v>0$满足$\pi_!v\neq v$（成长方向在不可复制维度上有正分量），则$\rho(U_{22}\circ\pi_!)<\rho(U_{22})$。
\end{enumerate}
\end{theorem}

\begin{proof}
(1) 由谱半径次可乘性，$\rho(U_{22}\pi_!)\leq\|U_{22}\pi_!\|\leq\|U_{22}\|\|\pi_!\|\leq\|U_{22}\|$。取算子范数为谱范数时$\rho(A)\leq\|A\|$，但$\rho(AB)\leq\rho(A)\|B\|$不直接给出$\rho(U_{22}\pi_!)\leq\rho(U_{22})$。更精确地：对任意$\epsilon>0$，存在算子范数$\|\cdot\|_\epsilon$使$\|U_{22}\|_\epsilon\leq\rho(U_{22})+\epsilon$（Yamaguchi定理/谱半径的范数刻画），故$\rho(U_{22}\pi_!)\leq(\rho(U_{22})+\epsilon)\|\pi_!\|_\epsilon\leq\rho(U_{22})+\epsilon$，令$\epsilon\to0$得证。

(2) 设$r=\rho(U_{22})$，由Perron--Frobenius定理，不可约非负矩阵$U_{22}$有唯一（相差正倍数）正特征向量$v>0$使$U_{22}v=rv$。假设$\rho(U_{22}\pi_!)=r$，则存在$x\geq0,x\neq0$使$U_{22}\pi_!x=rx$。由$U_{22}$不可约且$\pi_!x\leq x$（投影的单调性），$U_{22}\pi_!x\leq U_{22}x$。若$\pi_!x\neq x$，不可约性保证$U_{22}\pi_!x<U_{22}x$在某个正分量上严格成立，与$U_{22}\pi_!x=rx=U_{22}(cv)$矛盾（Perron向量唯一）。故$\pi_!x=x$，从而$U_{22}x=rx$，$x=cv$，因此$\pi_!v=v$，与条件矛盾。故严格不等。\qed
\end{proof}

\begin{remark}[生命论语境]\lvC
$\pi_!$丢失的是线性状态中不可被!-模态复制的部分——即交互维度（定理\ref{thm:dichotomy}：$X$在$\lolli$后件时阻止!-分配）。当生命的成长方向（Perron向量$v$）依赖于交互维度时（$v$在该维度有正分量），异化必然严格降低$N$。$N$从$>1$（创造/盈余）降到$=1$（持平）甚至$<1$（衰退/慢性死亡），与定理\ref{thm:slow-death}（慢性死亡）衔接。
\end{remark}

\begin{theorem}[革命=消除!-寄生Cut]\label{thm:revolution-cut}\lvI
消除!-寄生反馈回路（移除$\pi_!$和!-萃取$!_B$）后，$U_{22}$恢复为原始线性反馈，$N$恢复到$\geq1$，$W$降到0，$M$从有限恢复为无限。社会革命（切断资本的!-回路）和个人明性（切断意识形态的!-回路）在GoI中是同一个操作。
\end{theorem}

\begin{remark}[过渡时期]
Cut消除需要多步计算。过渡时期对应!-反馈已断但线性反馈未完全建立，$N$在波动中上升。系统保持生产性当且仅当线性反馈的增长速度超过!-残余的衰减速度。
\end{remark}

\section{综合}

\begin{theorem}[$G^\omega$是自函子]\label{thm:gomega-endofunctor}\lvP
$G^\omega:\mathbf{PTC}\to\mathbf{PTC}$是PTC范畴上的自函子，把PTC映射为dagger紧闭合PTC。由定理\ref{thm:f3-fixedpoint}，三层封闭：$G^{\omega 3}(\mathcal{C})\cong G^{\omega 2}(\mathcal{C})$。
\end{theorem}

\begin{proof}
\textbf{$G^\omega(\mathcal{C})$是PTC。}$G^\omega(\mathcal{C})$的态射是$\mathcal{C}$中的生产性Mealy机，复合用$\mathrm{Tr}^\omega$。逐条验证PTC公理：
\begin{itemize}
  \item \emph{Yanking}：$G^\omega(\mathcal{C})$中的弯线是$\mathcal{C}$中Mealy机的弯线，其行为态射由$\mathcal{C}$的Yanking公理保证为恒等流。
  \item \emph{Tightening}：$\mathrm{beh}((h\otimes\mathrm{id})f(g\otimes\mathrm{id}),x_0)=h^\omega\,\mathrm{beh}(f,x_0)\,g^\omega$由$\mathcal{C}$的Tightening公理直接推出（前后处理都是$\mathcal{C}$中的态射）。
  \item \emph{组合}：$\mathrm{beh}(f;g,(x_0,y_0))=\mathrm{beh}(g,y_0)\circ\mathrm{beh}(f,x_0)$由$\mathcal{C}$的组合公理推出（Mealy机复合在$\mathcal{C}$中用$\mathrm{Tr}^\omega$定义）。
  \item \emph{Productivity}：$\mathcal{C}$中生产性Mealy机的复合仍生产性（$\mathcal{C}$的Productivity公理），故$G^\omega(\mathcal{C})$中态射完全。
\end{itemize}

\textbf{dagger紧闭合。}定理\ref{thm:gomega-compact}已证$G^\omega(\mathcal{C})$是dagger compact closed范畴（对偶$(A^+,A^-)^*=(A^-,A^+)$，蛇方程由Yanking保证，类比Joyal--Street--Verity 1996的Int-构造）。

\textbf{函子性。}对PTC态射$F:\mathcal{C}\to\mathcal{D}$（保持张量、迹和生产性的强幺半函子），定义$G^\omega(F):G^\omega(\mathcal{C})\to G^\omega(\mathcal{D})$逐对象作用$F$、逐态射将$\mathcal{C}$中Mealy机$(X,R,x_0)$映射为$(FX,F(R),F(x_0))$。$F$保持$\mathrm{Tr}^\omega$（PTC态射定义），故$G^\omega(F)$保持复合和恒等。$G^\omega(\mathrm{id}_\mathcal{C})=\mathrm{id}_{G^\omega(\mathcal{C})}$，$G^\omega(F\circ G)=G^\omega(F)\circ G^\omega(G)$。

\textbf{三层封闭。}由定理\ref{thm:f3-fixedpoint}，$f^3$是层级提升的不动点。$G^\omega$构造对应增加一层反馈嵌套（$G^{\omega(n+1)}(\mathcal{C})$的态射比$G^{\omega(n)}(\mathcal{C})$多一条弯线），故$G^{\omega 3}(\mathcal{C})\cong G^{\omega 2}(\mathcal{C})$。\qed
\end{proof}

\begin{theorem}[生命论基本定理]\label{thm:fundamental}\lvI
以下九条命题共享同一数学结构——\textbf{线性资源不可被$\bang$-模态还原}——其中四条有严格证明、两条分拆为"数学+等同"、三条为纯诠释：
\begin{enumerate}
\item 生命是自指过程$S=f(S)$——[诠释]（践演元预设；数学载体：定义7与GoI执行语义）
\item 生命过程是线性的（资源敏感，不可复制）——[定理]（公理2）
\item 生命流$\nu F_2$不可被$\bang$-资本化——[定理]（定理\ref{thm:no-bang-stream}、\ref{thm:bang-gomega}；种子定理\ref{thm:no-A-to-bangA}的GoI推广）
\item "运行被沉积模型接管"是[定理]（定理\ref{thm:alienation-mistake}）；与"异化"的等同是[诠释]
\item "维持运行与沉积模型的区分"是[定理]（定理\ref{thm:mingxing}）；与"明性"的等同是[诠释]
\item "消除$\bang$-寄生Cut"的GoI操作是[定理]（Cut消除）；与"革命"的等同是[诠释]
\item 量子不可克隆是同一结构在物理领域的显现——[定理]（定理\ref{thm:no-cloning-corollary}）
\item "$N<1$时$M$有限"的几何级数是[定理]（定理\ref{thm:slow-death}）；"$M$=意义、慢性死亡"是[诠释]
\item 生命过程超越有限证明——[定理]（定理\ref{thm:life-beyond-proof}）
\end{enumerate}

\textbf{此处"共享同一结构"的精确含义}：每条命题都可从种子定理\ref{thm:no-A-to-bangA}（不存在$A\to\bang A$）或其经不动点/GoI结构的推广（定理\ref{thm:dichotomy}、\ref{thm:no-bang-stream}、\ref{thm:bang-gomega}）与线性逻辑公理的一个小推导中得到；第1、6、8条虽为[诠释]，其数学载体同样来自定理\ref{thm:no-A-to-bangA}/\ref{thm:dichotomy}/\ref{thm:no-bang-stream}/\ref{thm:bang-gomega}。

本文不声称"从践演事实必然推出全部九条"。声称的是：各条共享上述种子结构；其中2、3、7、9有严格证明（[定理]）；4、5、6的数学载体是[定理]而哲学等同是[诠释]；1、8为[诠释]（数学部分分别在践演预设与注22）。
\end{theorem}

\section{相关工作与公理精简}

\begin{remark}[与Guarded Recursion的关系]
PTC的Productivity公理与Nakano（2000）later模态$\rhd$的guardedness条件是同一概念：流对象$A^\omega=\nu X.(A\tensor X)$对应$\mathrm{Stream}\ A=A\times\rhd(\mathrm{Stream}\ A)$，生产性迹对应guarded fixed point，Mealy机对应因果流函数。PTC的新贡献是$\bang$与生产性的交互（定理\ref{thm:bang-no-ptc}）：$\bang$不穿透guarded/生产性状态——这在guarded recursion文献中未被研究。
\end{remark}

\begin{theorem}[Composition是定理不是公理]\label{thm:composition-derivable}\lvP
PTC的Composition公理（生产性Mealy机反馈复合仍生产性）可从Tightening和Yanking推出，不是独立公理。PTC的独立公理为三条：SMC+流对象、迹公理（Yanking+Tightening）、Productivity。
\end{theorem}

\begin{proof}
在traced monoidal category中，复合的结合律从tightening/naturality和yanking推出（Joyal-Street-Verity 1996）。PTC中Mealy机复合用生产性迹定义，其结合律和同一律由相同论证给出。Productivity独立于标准迹公理：标准迹允许nilpotency（$U_{22}^n=0$），但不要求无限反馈每步有产出。
\end{proof}

\begin{remark}
定理\ref{thm:composition-derivable}将PTC公理精简为三条（SMC+流对象、迹公理、Productivity），与公理3的Seely条件无冲突：Seely是$\bang$模态的额外结构，不属于迹范畴公理。二者分别在§2与§18独立引入。
\end{remark}

\section{待证猜想}

\begin{conjecture}[生命测度]\label{conj:measure}
操作范畴上可定义一个测度，使量论公式$M=\alpha\cdot T\cdot N$成为定量定理，其中$T$是存活区间测度、$N$是余代数展开的净变化方向、$\alpha$是f-层级系数。目前$M$仅在加权$\mathbf{Rel}$中有独立于哲学诠释的纯数学定义（§20"模型内定义"）。
\end{conjecture}

\begin{conjecture}[f\textsuperscript{3}异化定理]\label{conj:f3}
\textbf{已修正表述。}原猜想"$f^1$和$f^2$层不存在到$\bang$-模态自我模型的态射"不成立——$f^2$层有$\bang\ulcorner f\urcorner$（定理\ref{thm:self-model-bang}）。正确的区分是：$f^2$可被\emph{控制}（替换模型后直接运行），但不能被\emph{异化}（异化需要"自愿认同"的选择能力，只有$f^3$具备）。命题\ref{prop:f3-alienation}已按此修正。严格化待办：在GoI反馈线中形式化"判断步骤"（$f^3$的第三条反馈线如何实现"选择是否按模型运行"），使"控制vs异化"的区分获得数学表达。
\end{conjecture}

\begin{conjecture}[践演判断的一致性]\label{conj:performative-consistent}
践演判断$\triangleright A$可一致地加入Martin-L\"of类型论而不改变证明论强度。
\end{conjecture}

\begin{conjecture}[动态自毁]\label{conj:dynamic}
定理\ref{thm:parasite-self-destruct}的动态版本：在操作范畴的traced/timed扩展中，反自指寄生体的持续作用必然使$S$趋向$\I$。
\end{conjecture}

\begin{conjecture}[$\bang$不穿透PTC的完全一般性]\label{conj:bang-ptc}
\textbf{已由定理\ref{thm:bang-no-ptc}解决。}在任意操作范畴中，$F_2$型Mealy机不存在保持余归纳迹的自然!-提升。证明使用三个独立论证：(1)$\kappa$-分叉破坏单一线索；(2)$\eps$不提供唯一状态选择；(3)$\bang$与$\nu$不交换使类型不匹配。三个论证均只依赖公理1--4，不依赖特定模型。
\end{conjecture}

\vspace{1em}
\noindent\textbf{猜想列表 $\leftrightarrow$ [条件推论] 对应表：}
\begin{center}
\begin{tabular}{ll}
\hline
\textbf{[条件推论]} & \textbf{状态} \\
\hline
定理\ref{thm:growth-needs-linear}（增长须消耗线性资源） & \textbf{已升级为定理}：Seely同构(公理3iii)强制δ对角，!-模态态射只能复制/重排/丢弃，与定理6对偶（!G不能产生线性A） \\
定理\ref{thm:godel-f2}（f$^2$哥德尔困境） & \textbf{已严格化}：f²转移函数图灵可计算→PA可表示；!⌜f⌝=哥德尔码；对角线引理给出G；f³第三条反馈线对应元层面，▷G对应系统外操作 \\
定理\ref{thm:f3-fixedpoint}（f$^3$不动点） & \textbf{已升级为[定理]}：编码论证证明$H_3$穷尽所有自指变换 \\
定理\ref{thm:extraction-alienation}（萃取加速异化） & \textbf{已升级为定理}：一般PTC层面由定理bang-no-ptc的κ-分叉论证覆盖；加权Rel中N(W)=(1-W)N(0) \\
定理\ref{thm:coh-ptc}（Coh$^\omega$是PTC） & \textbf{已升级为[定理]}：clique归纳证明+PTC四公理逐条验证 \\
定理\ref{thm:alienation-lowers-n}（异化降低$N$） & \textbf{加权$\mathbf{Rel}$中已严格证明}：Perron--Frobenius给出严格下降条件；谱半径论证本质依赖非负矩阵，为模型相关定理 \\
定理\ref{thm:gomega-endofunctor}（$G^\omega$是自函子） & \textbf{已升级为[定理]}：PTC四公理从底范畴继承+dagger紧闭合+函子性+三层封闭 \\
定理\ref{thm:bang-no-ptc}（$\bang$不穿透生产性迹） & \textbf{已升级为[定理]}：三论证（κ分叉/ε不唯一/!≇ν!ν）纯范畴证明 \\
\hline
\end{tabular}
\end{center}

\section{总结}

本文定义了操作范畴，证明了56个定理，核心结构如下：

\begin{center}
\begin{tabular}{lll}
\hline
\textbf{数学结构} & \textbf{生命论} & \textbf{马克思} \\
\hline
线性对象$A$ & 阳/活操作 & 活劳动 \\
$\bang$-模态$\bang A$ & 阴/沉积 & 死劳动/资本 \\
收缩$\bang A\to\bang A\tensor\bang A$ & 阴的复制 & 资本增殖 \\
弃置$\bang A\to A$ & 阴激活为阳 & 资本投入生产 \\
无$A\to\bang A$（定理\ref{thm:no-A-to-bangA}） & 阳不可自动阴化 & 劳动不能自动变资本 \\
无$S\to\bang S$（定理\ref{thm:life-not-capital}） & 生命不可完全资本化 & 异化有数学限度 \\
共Kleisli范畴（CCC） & 阴的领域 & 流通领域/商品拜物教 \\
基范畴（SMCC） & 阳的领域 & 生产领域/剥削可见 \\
$\mN F$（最终余代数） & 生命自指 & 活劳动/生命流 \\
$\mU F$（初始代数） & 沉积自指 & 资本积累 \\
$\bang\nu F_1\to\nu(F_1\bang)$（定理\ref{thm:nu-lift-data}） & 数据流可复制 & 录像/文本可复制 \\
$\bang\nu F_2\not\to\nu(F_2\bang)$（定理\ref{thm:nu-no-lift-life}） & 生命流不可复制 & 人不可复制 \\
$\bang\mu F_2\not\to\mu(F_2\bang)$（定理\ref{thm:mu-no-lift}） & 有限交互不可复制为过程 & 对话记录不是对话 \\
迹$\mathrm{Tr}(f)$（命题\ref{prop:mealy-goi}） & 藏起自我只留行为 & 测量/实际从属 \\
余归纳执行（定义\ref{def:coinductive-execution}） & 生命每步有产出 & 活出完整 \\
\hline
\end{tabular}
\end{center}

\vspace{1em}
\noindent\textbf{结构说明。}表中所有"数学结构"列的条目均为[定理]或[结构]级；"生命论""马克思"两列为[诠释]级对应。本文的数学贡献集中在"数学结构"列——尤以定理\ref{thm:no-A-to-bangA}（种子）、定理\ref{thm:dichotomy}（二分法）、定理\ref{thm:no-bang-stream}/\ref{thm:bang-no-ptc}/\ref{thm:bang-gomega}（$\bang$不穿透流处理器）、定理\ref{thm:no-cloning-corollary}（不可克隆为推论）、定理\ref{thm:life-beyond-proof}（超越有限证明）为主线；两侧为哲学赋值的[诠释]，不声称数学蕴含。

\vspace{1em}
\noindent\textbf{致谢。} 本文的哲学框架来自生命论（明本论），其核心命题——操作先于实体、自指维持、反自指寄生、践演坐实——均由生命论提出。数学部分建立在线性逻辑、范畴论与马克思政治经济学的交汇之上。

\begin{thebibliography}{99}
\bibitem{lawvere69} F.W.~Lawvere. Diagonal arguments and cartesian closed categories. \textit{Category Theory, Homology Theory and their Applications II}, 1969.
\bibitem{girard87} J.-Y.~Girard. Linear logic. \textit{Theoretical Computer Science}, 50(1):1--102, 1987.
\bibitem{grishin82} V.N.~Grishin. A nonstandard logic and its application to set theory. \textit{Studies in Formalized Languages and Nonclassical Logics}, 1982.
\bibitem{aczel88} P.~Aczel. \textit{Non-Well-Founded Sets}. CSLI, 1988.
\bibitem{benton95} P.N.~Benton. A mixed linear and non-linear logic: Proofs, terms and models. \textit{CSL '94}, LNCS 933, 1995.
\bibitem{baelde07} D.~Baelde and D.~Miller. Least and greatest fixed points in linear logic. \textit{LPAR}, 2007.
\bibitem{roberts23} D.M.~Roberts. Substructural fixed-point theorems and the diagonal argument. \textit{Compositionality}, 5(8), 2023.
\bibitem{martinlof84} P.~Martin-L\"of. \textit{Intuitionistic Type Theory}. Bibliopolis, 1984.
\bibitem{jacobs05} B.~Jacobs. \textit{Introduction to Coalgebra}. 2005.
\bibitem{marx67} K.~Marx. \textit{Das Kapital}, Band I. 1867.
\bibitem{girard89} J.-Y.~Girard. Geometry of Interaction I: Interpretation of system F. \textit{Proc. Logic Colloquium '88}, 1989.
\bibitem{abramsky96} S.~Abramsky. Retracing some paths in process algebra. \textit{CONCUR '96}, LNCS 1119, 1996.
\bibitem{jsv96} A.~Joyal, R.~Street, D.~Verity. Traced monoidal categories. \textit{Math. Proc. Camb. Phil. Soc.}, 119:447--468, 1996.
\bibitem{hasegawa16} M.~Hasegawa. Memoryful Geometry of Interaction. \textit{POPL 2016}.
\bibitem{braverman74} H.~Braverman. \textit{Labor and Monopoly Capital}. 1974.
\bibitem{nakano00} H.~Nakano. A modality for recursion. \textit{LICS 2000}.
\bibitem{abramsky04} S.~Abramsky and B.~Coecke. A categorical semantics of quantum protocols. \textit{LICS 2004}.
\bibitem{wootters82} W.K.~Wootters and W.H.~Zurek. A single quantum cannot be cloned. \textit{Nature}, 299:802--803, 1982.
\bibitem{mogelberg14} R.E.~Møgelberg. A type theory for productive coprogramming via guarded recursion. \textit{CSL-LICS 2014}.
\bibitem{danos95} V.~Danos and L.~Regnier. Proof-nets and the Hilbert space. \textit{Advances in Linear Logic}, 1995.
\end{thebibliography}

\end{document}
